\documentclass[11pt,oneside]{article}

\usepackage[T1]{fontenc}
\usepackage{lmodern}
\usepackage{microtype}
\usepackage{geometry}
\geometry{margin=1in,headheight=14pt}
\usepackage{amsmath,amssymb,amsthm,mathtools}
\usepackage{booktabs,longtable,tabularx,array,multirow}
\usepackage{enumitem}
\usepackage{xcolor}
\usepackage{xurl}
\usepackage{hyperref}
\usepackage{listings}
\usepackage{fancyhdr}
\usepackage{titlesec}
\usepackage{needspace}

\definecolor{proofgreen}{RGB}{15,105,68}
\definecolor{inheritedblue}{RGB}{24,91,145}
\definecolor{openviolet}{RGB}{98,55,145}
\definecolor{cautionred}{RGB}{150,55,45}
\definecolor{codegray}{RGB}{248,248,248}
\definecolor{rulegray}{RGB}{92,101,111}

\hypersetup{
  colorlinks=true,
  linkcolor=inheritedblue,
  citecolor=proofgreen,
  urlcolor=inheritedblue,
  pdftitle={Alaoglu--Erdos Continuation Report 28: Confluent Functional Calculus and Rational Self-Extensions},
  pdfauthor={OpenAI research continuation},
  pdfsubject={New exact partial results and a sharpened research program for the Alaoglu--Erdos integer-exponent conjecture}
}
\urlstyle{same}

\pagestyle{fancy}
\fancyhf{}
\lhead{Alaoglu--Erd\H{o}s continuation report 28}
\rhead{Confluence and rational self-extensions}
\cfoot{\thepage}
\renewcommand{\headrulewidth}{0.4pt}

\titleformat{\section}{\Large\bfseries}{\thesection}{0.75em}{}
\titleformat{\subsection}{\large\bfseries}{\thesubsection}{0.75em}{}
\titleformat{\subsubsection}{\normalsize\bfseries}{\thesubsubsection}{0.75em}{}

\setlist[itemize]{leftmargin=2em,itemsep=0.28em,topsep=0.4em}
\setlist[enumerate]{leftmargin=2.2em,itemsep=0.28em,topsep=0.4em}
\setlist[description]{leftmargin=3.5cm,labelsep=0.55cm,itemsep=0.35em,topsep=0.4em}
\setlength{\parindent}{0pt}
\setlength{\parskip}{0.55em}
\setlength{\emergencystretch}{3.5em}

\newtheorem{theorem}{Theorem}[section]
\newtheorem{proposition}[theorem]{Proposition}
\newtheorem{lemma}[theorem]{Lemma}
\newtheorem{corollary}[theorem]{Corollary}
\newtheorem{conjecture}[theorem]{Conjecture}
\newtheorem{criterion}[theorem]{Closing criterion}
\newtheorem{question}[theorem]{Question}
\theoremstyle{definition}
\newtheorem{definition}[theorem]{Definition}
\newtheorem{construction}[theorem]{Construction}
\newtheorem{target}[theorem]{Research target}
\newtheorem{audit}[theorem]{Audit finding}
\theoremstyle{remark}
\newtheorem{remark}[theorem]{Remark}
\newtheorem{warning}[theorem]{Caution}

\newcommand{\N}{\mathbb N}
\newcommand{\Nzero}{\mathbb N_0}
\newcommand{\Z}{\mathbb Z}
\newcommand{\Q}{\mathbb Q}
\newcommand{\R}{\mathbb R}
\newcommand{\C}{\mathbb C}
\newcommand{\Qbar}{\overline{\mathbb Q}}
\newcommand{\Mat}{\operatorname{M}}
\newcommand{\End}{\operatorname{End}}
\newcommand{\tr}{\operatorname{tr}}
\newcommand{\im}{\operatorname{im}}
\newcommand{\ad}{\operatorname{ad}}
\newcommand{\diag}{\operatorname{diag}}
\newcommand{\Sym}{\operatorname{Sym}}
\newcommand{\Ext}{\operatorname{Ext}}
\newcommand{\den}{\operatorname{den}}
\newcommand{\lcm}{\operatorname{lcm}}
\newcommand{\sgn}{\operatorname{sgn}}
\newcommand{\Span}{\operatorname{span}}
\newcommand{\cE}{\mathcal E}
\newcommand{\cC}{\mathcal C}
\newcommand{\cQ}{\mathcal Q}
\newcommand{\cR}{\mathcal R}
\newcommand{\cZ}{\mathcal Z}
\newcommand{\cO}{\mathcal O}


\newcommand{\dd}{\mathop{}\!\mathrm d}
\newcommand{\AEname}{Alaoglu--Erd\H{o}s}
\newcommand{\paperproof}{\textcolor{proofgreen}{\textbf{[proved here]}}}
\newcommand{\inherited}{\textcolor{inheritedblue}{\textbf{[inherited]}}}
\newcommand{\openstatus}{\textcolor{openviolet}{\textbf{[open interface]}}}
\newcommand{\cautionstatus}{\textcolor{cautionred}{\textbf{[not a proof of the conjecture]}}}
\newcommand{\lean}[1]{\texttt{\detokenize{#1}}}

\newenvironment{statusbox}[1]
{\begin{center}\begin{minipage}{0.94\textwidth}\raggedright\hbadness=10000\hrule
\vspace{0.45em}\textbf{Status.} #1\par\vspace{0.35em}}
{\vspace{0.35em}\hrule\end{minipage}\end{center}}

\newenvironment{resultbox}
{\begin{center}\begin{minipage}{0.94\textwidth}\hrule\vspace{0.45em}}
{\vspace{0.35em}\hrule\end{minipage}\end{center}}

\lstdefinestyle{leanstyle}{
  basicstyle=\ttfamily\footnotesize,
  backgroundcolor=\color{codegray},
  frame=single,
  rulecolor=\color{rulegray},
  showstringspaces=false,
  breaklines=true,
  columns=fullflexible,
  keepspaces=true,
  tabsize=2,
  captionpos=b
}
\lstdefinestyle{pythonstyle}{
  language=Python,
  basicstyle=\ttfamily\footnotesize,
  backgroundcolor=\color{codegray},
  frame=single,
  rulecolor=\color{rulegray},
  numbers=left,
  numberstyle=\tiny\color{rulegray},
  numbersep=8pt,
  showstringspaces=false,
  breaklines=true,
  columns=fullflexible,
  keepspaces=true,
  tabsize=4,
  captionpos=b
}

\begin{document}
\pagenumbering{gobble}

\begin{titlepage}
\centering
\vspace*{0.55cm}
{\Huge\bfseries The Alaoglu--Erd\H{o}s\par}
\vspace{0.12cm}
{\Huge\bfseries Integer-Exponent Conjecture\par}
\vspace{0.65cm}
{\LARGE Continuation Report 28\par}
\vspace{0.18cm}
{\Large\bfseries Confluent Functional Calculus, Rational Self-Extensions,\par}
\vspace{0.08cm}
{\Large\bfseries and the Missing Centralizer Jet\par}
\vspace{0.8cm}
{\large New additive progress beyond the unified 27-report manuscript\par}
\vspace{0.22cm}
{\large August 22, 2026\par}
\vfill
\begin{minipage}{0.91\textwidth}
\small
\textbf{Bottom line.}
The conjecture remains unproved in this report.  The main advance is instead an exact
classification theorem for every first-order rational self-extension of a semisimple matrix
whose spectral values are already known.  For a hypothetical nonintegral exponent $x$, the
$x$-th power of such an extension is rational \emph{exactly} in the split, commutator
directions.  Every genuinely non-split rational extension exposes the missing scalar $x$.
Equivalently, real power acts on the rational self-extension quotient as $x$ times a rational
automorphism.  This turns the informal slogan ``a Jordan block would solve the problem'' into
a complete rationality-locus theorem.

A second advance develops a confluent approximation hierarchy.  Symmetric transported secants
improve the first-order error in the unified report from $O(h)$ to $O(h^2)$, and explicit
Richardson coefficients raise this to $O(h^{2m})$ for arbitrary fixed $m$, while all
approximants remain rational under the counterexample hypothesis.  A denominator audit shows
precisely why this still does not force rationality: the available heights are exponential in
the Diophantine-approximation parameter, whereas the analytic gain is polynomial.  The report
also proves a toric-functor no-go theorem, gives a genuinely integral $3\times3$ Newton-chain
example, and isolates an affine collision program based on the additive relation $3=2+1$.
\end{minipage}
\vfill
{\small Prepared as a standalone continuation; no theorem below is represented as Lean-checked
unless explicitly stated otherwise.\par}
\end{titlepage}

\clearpage
\pagenumbering{roman}
\begin{abstract}
Assume that $x\in\R$ and
\[
 M:=2^x\in\Z,\qquad N:=3^x\in\Z.
\]
The attached unified manuscript establishes a large body of exact reductions, finite
verification, determinant estimates, and conditional structure, while leaving the original
conjecture open.  This continuation deliberately avoids repeating that material.  Its first
purpose is to identify, in a basis-free way, the exact finite-dimensional datum absent from all
semisimple constructions.

Let $A\in\Mat_n(\Q)$ be semisimple with positive rational eigenvalues
$\lambda_1,\dots,\lambda_s$, rational spectral projectors, and
$\lambda_i^x\in\Q$.  For $H\in\Mat_n(\Q)$ form the first-order block extension
\[
 \cE_A(H)=\begin{pmatrix}A&H\\0&A\end{pmatrix}.
\]
The analytic block functional calculus gives
\[
 \cE_A(H)^x=
 \begin{pmatrix}A^x&D(z^x)_A[H]\\0&A^x\end{pmatrix}.
\]
If $x\notin\Q$, then
\[
 D(z^x)_A[H]\in\Mat_n(\Q)
 \quad\Longleftrightarrow\quad
 H\in[A,\Mat_n(\Q)].
\]
Thus rationality survives exactly on extensions already split over $\Q$; every nonzero class
in $\Mat_n(\Q)/[A,\Mat_n(\Q)]$ detects $x$.  On that quotient the derivative is
$x$ times multiplication by the rational invertible matrix $A^{x-1}$.  Several trace and
matrix-coefficient extraction criteria follow.

The second purpose is to connect this exact obstruction with rational confluence.  For
$r=2^a3^b$ and $r^x=M^aN^b$, define
\[
 J_x(r)=\frac{r^x-r^{-x}}{r-r^{-1}}
       =\frac{\sinh(xh)}{\sinh h},\qquad h=\log r.
\]
Then $J_x(r)\in\Q$ and $J_x(e^h)=x+O_x(h^2)$.  With
\[
 c_{m,k}=(-1)^{k-1}\frac{2(m!)^2}{(m-k)!(m+k)!},
 \qquad
 R_{m,x}(r)=\sum_{k=1}^m c_{m,k}J_x(r^k),
\]
we prove $R_{m,x}(e^h)=x+O_{x,m}(h^{2m})$.  This is a genuine higher-order strengthening of
the transported secants already present in the unified report, but the reduced-denominator
problem remains decisive.
\end{abstract}

\tableofcontents
\clearpage
\pagenumbering{arabic}

\section{Executive answer and evidentiary boundary}
\label{sec:executive}

\subsection{Outcome}

The requested full proof was not obtained.  What was obtained is a collection of exact new
results that sharply narrows one previously qualitative obstruction:

\begin{enumerate}
\item \textbf{Complete first-order rationality locus.}
For every rational semisimple base $A$ with rational positive spectrum and rational powered
spectrum, a hypothetical irrational exponent makes the rational block extensions with rational
$x$-th power \emph{exactly} the commutator, hence split, extensions.

\item \textbf{Extension-space action.}
Modulo rational splitting, the Fr\'echet derivative of $z^x$ is exactly $x$ times a rational
automorphism.  This is the finite-dimensional location of the missing arithmetic datum.

\item \textbf{Toric no-go.}
Direct sums, tensor products, duals, symmetric powers, exterior powers, Schur functors, and
rational subquotients of the split torus representation generated by $2$ and $3$ can create
multiplicity but never a nonzero extension class.  Spectral collision is not confluence.

\item \textbf{Integral collision examples.}
An explicit integral $3\times3$ Newton chain has integral $x$-th power under the original
hypotheses, is not unimodularly diagonalizable, and produces repeated eigenvalues after a
tensor square; nevertheless all collisions remain semisimple.

\item \textbf{Higher-order rational confluence.}
A symmetric secant and its Richardson extrapolants yield rational approximations to $x$ of
arbitrary fixed even order.  The exact denominator growth that blocks unit locking is isolated.

\item \textbf{Affine escape route.}
The family $\left(\begin{smallmatrix}q&1\\0&3\end{smallmatrix}\right)$, sampled at
$3$-smooth rational $q$, converges to the decisive Jordan extension at $q=3$.  Proving an
arithmetic compactness statement for these rational secants would close the conjecture.
\end{enumerate}

\begin{statusbox}{The classification, quotient, toric no-go, explicit matrix formulas,
degeneration theorem, Richardson identities, and conditional locking statements are
\paperproof{} paper arguments in this report.  They have been checked algebraically and by
symbolic scripts reproduced in Appendix~\ref{app:symbolic}, but they are not claimed to have
been accepted by a proof assistant.  The original \AEname{} conjecture remains \openstatus{}.}
\end{statusbox}

\subsection{External reconnaissance as of August 22, 2026}
\label{sec:recon}

Because the prompt specifically referred to recent public AI-assisted mathematics, the public
record was checked separately from the mathematical derivation.  Anthropic publicly launched
Claude Fable~5 on June 9, 2026 and restored broad access on July 1 after a temporary
suspension~\cite{AnthropicFable,AnthropicRedeploy}.  The recent Jacobian development is a
disproof in dimensions greater than two, with explicit constructions and exact verification;
it does not settle the still-distinct two-dimensional case~\cite{GaoJacobian}.  A public
MathOverflow update and the ICARM leaderboard report a curve of Mordell--Weil rank at least
$30$, with exact rank $30$ conditional on BSD and GRH~\cite{SilvermanMO,ICARMRank}.  Public
computational announcements also report exact constructions for the twelve previously open
Hadamard orders below $2000$~\cite{VibeMathedHadamard,OEISHadamard}; at the time of writing this
is best treated as an exact-computation announcement rather than a conventional refereed
paper.

No public source located in this search announced a proof of the \AEname{} conjecture or a Lean
formalization of one.  The rival team's reported result may therefore be private or
unannounced.  Nothing in this report assumes access to it.

\subsection{The immovable transcendence boundary}

The unified manuscript already explains that the conjecture is an integral special case of the
four exponentials problem.  The recent Matrix Coefficient Conjecture of Dasgupta and
Kakde gives a particularly clean modern language for this boundary; its $2\times2$ case is
equivalent to four exponentials~\cite{DasguptaKakde}.  No claim below silently assumes that
conjecture.  The point of the present work is instead to seek extra arithmetic structure coming
from integrality, lattices, and confluence.

\section{Non-duplication ledger}
\label{sec:ledger}

The attached manuscript is nearly one megabyte and already contains twenty-seven
continuation reports.  Repeating its reductions would obscure rather than advance the project.
Table~\ref{tab:ledger} records what is inherited and what is genuinely added here.

\begin{longtable}{>{\raggedright\arraybackslash}p{0.25\textwidth}
                  >{\raggedright\arraybackslash}p{0.32\textwidth}
                  >{\raggedright\arraybackslash}p{0.34\textwidth}}
\caption{Non-duplication ledger.}\label{tab:ledger}\\
\toprule
Topic & Already present in the unified manuscript & New contribution in Report 28 \\
\midrule
\endfirsthead
\toprule
Topic & Already present in the unified manuscript & New contribution in Report 28 \\
\midrule
\endhead
Basic reductions & Positivity, rational descent, transcendence of a hypothetical
counterexample, exact conditional semigroup structure, finite verification & Used only as
black-box baseline; not reproved except for two one-line closure lemmas \\
\addlinespace
Four exponentials & Exact reduction and Matrix Coefficient formulation & No new transcendence
claim; boundary is kept explicit \\
\addlinespace
Semisimple matrices & Formula for a matrix satisfying $(T-2I)(T-3I)=0$ and the observation that
distinct-node functional calculus is blind & Full Fr\'echet derivative decomposition for every
rational semisimple $A$, with an if-and-only-if rationality locus \\
\addlinespace
Jordan blocks & A sufficient criterion from $(bI+N)^x=b^xI+xb^{x-1}N$ & Classification of all
first-order rational self-extensions, the extension quotient, and basis-free extraction maps \\
\addlinespace
Tensor/symmetric constructions & Qualitative semisimple blindness & Categorical torus
no-go theorem: multiplicity without extension, including rational subquotients \\
\addlinespace
Transported secants & Rational one-sided secants with first-order error and finite-jet ceilings &
Symmetric $O(h^2)$ secants and explicit Richardson extrapolation of order $2m$ \\
\addlinespace
Integral interpolation & Geometric Newton identities and Smith-form analyses & Explicit
$3\times3$ lattice-glued chain at eigenvalues $2,3,4$, plus tensor collision and
unbounded-conjugator degeneration theorems \\
\addlinespace
Operator/affine ideas & External affine ruler as a sufficient condition & A parameterized affine
collision at $q=3$ sampled entirely on the $3$-smooth rational group \\
\addlinespace
Lean plan & Many accepted modules and a prior Jordan closing lemma & A modular formalization
plan for the new algebraic classification and Richardson identities \\
\bottomrule
\end{longtable}

\section{Baseline, notation, and the exact question}
\label{sec:baseline}

Assume throughout that
\begin{equation}
 x\in\R,\qquad M=2^x\in\Z,\qquad N=3^x\in\Z.
 \label{eq:baseline}
\end{equation}
If $x<0$, then $0<2^x<1$, impossible for an integer; hence $x\ge0$.  The case $x=0$ is
integral.  In any hypothetical counterexample, therefore,
\begin{equation}
 x>0,
 \qquad x\notin\Q,
 \qquad M,N\ge2.
 \label{eq:hypothetical}
\end{equation}
The irrationality assertion is inherited: if $x=a/b\in\Q$ in lowest terms and $2^x\in\Q$,
unique factorization forces $b=1$.  The unified report further proves that a nonintegral
counterexample would be transcendental.

Let
\begin{equation}
 \Gamma=\{2^a3^b:a,b\in\Z\}\subset\Q_{>0}.
 \label{eq:gamma}
\end{equation}
Then \eqref{eq:baseline} gives an exact order-preserving group homomorphism
\begin{equation}
 \rho_x:\Gamma\longrightarrow\Q_{>0},
 \qquad
 \rho_x(2^a3^b)=M^aN^b=(2^a3^b)^x.
 \label{eq:rho}
\end{equation}
Because $\log2/\log3\notin\Q$, the logarithms of elements of $\Gamma$ form a dense additive
subgroup of $\R$.  In particular, there are sequences $r_j\in\Gamma\setminus\{1\}$ with
$r_j\to1$.  This density creates confluence analytically; the entire problem is to make that
confluence arithmetically rigid.

\begin{definition}[Divided differences]
For a function $f$ on positive reals, write
\[
 f[\lambda,\mu]=
 \begin{cases}
 \dfrac{f(\lambda)-f(\mu)}{\lambda-\mu},&\lambda\ne\mu,\\[1.1ex]
 f'(\lambda),&\lambda=\mu.
 \end{cases}
\]
For $f_x(z)=z^x$, the diagonal value is
$f_x[\lambda,\lambda]=x\lambda^{x-1}$.
\end{definition}

The distinction between the off-diagonal and diagonal divided differences is the organizing
principle of this report.  If $\lambda,\mu,\lambda^x,\mu^x\in\Q$ and $\lambda\ne\mu$, then
$f_x[\lambda,\mu]\in\Q$ automatically.  At a repeated node, rationality of the divided
difference is exactly rationality of $x\lambda^{x-1}$, hence---once $\lambda^x$ is known---of
$x$ itself.

\section{First-order block functional calculus}
\label{sec:blockfc}

\subsection{The universal block identity}

For $A,H\in\Mat_n(\C)$ define
\begin{equation}
 \cE_A(H)=
 \begin{pmatrix}
 A&H\\0&A
 \end{pmatrix}.
 \label{eq:blockextension}
\end{equation}
This is the universal first-order self-extension of $A$ in the direction $H$.

\begin{theorem}[Block derivative identity]
\label{thm:blockderivative}
Let $f$ be holomorphic on a neighborhood of the spectrum of $A$.  Then
\begin{equation}
 f(\cE_A(H))=
 \begin{pmatrix}
 f(A)&Df_A[H]\\0&f(A)
 \end{pmatrix},
 \label{eq:blockderivative}
\end{equation}
where $Df_A$ is the Fr\'echet derivative of the holomorphic functional calculus at $A$.
\end{theorem}

\begin{proof}
For a positively oriented contour $\gamma$ enclosing the spectrum,
\[
 f(B)=\frac1{2\pi i}\int_\gamma f(z)(zI-B)^{-1}\,\dd z.
\]
The block inverse is exact:
\[
 (zI-\cE_A(H))^{-1}=
 \begin{pmatrix}
 (zI-A)^{-1}&(zI-A)^{-1}H(zI-A)^{-1}\\
 0&(zI-A)^{-1}
 \end{pmatrix}.
\]
Integrating entrywise gives \eqref{eq:blockderivative}; the upper-right contour integral is
the standard Fr\'echet derivative.
\end{proof}

For the positive-spectrum matrices used below, $f_x(z)=z^x$ is taken on the principal branch
near the spectrum, so Theorem~\ref{thm:blockderivative} applies for every real $x$.

\subsection{Spectral divided-difference formula}

\begin{theorem}[Semisimple derivative]
\label{thm:semisimplederivative}
Let $A$ be semisimple with distinct eigenvalues $\lambda_1,\dots,\lambda_s$ and spectral
projectors $P_1,\dots,P_s$.  Then
\begin{equation}
 Df_A[H]=\sum_{i,j=1}^s f[\lambda_i,\lambda_j]P_iHP_j.
 \label{eq:divideddifference}
\end{equation}
In particular, for $f_x(z)=z^x$,
\begin{align}
 Df_{x,A}[H]
 &=\sum_{i\ne j}\frac{\lambda_i^x-\lambda_j^x}{\lambda_i-\lambda_j}P_iHP_j
   +x\sum_i\lambda_i^{x-1}P_iHP_i.
 \label{eq:powerderivative}
\end{align}
\end{theorem}

\begin{proof}
Insert $(zI-A)^{-1}=\sum_i(z-\lambda_i)^{-1}P_i$ into the contour formula in the proof of
Theorem~\ref{thm:blockderivative}.  For $i\ne j$, partial fractions give the divided
difference; for $i=j$, Cauchy's differentiation formula gives $f'(\lambda_i)$.
\end{proof}

Equation~\eqref{eq:powerderivative} is already suggestive.  Every off-diagonal coefficient is
a secant and is rational from value data.  Every diagonal block is multiplied by $x$.
The next section turns that observation into an exact if-and-only-if theorem.

\section{Exact rationality locus for first-order self-extensions}
\label{sec:rationalitylocus}

\subsection{Centralizer and commutator directions}

Let
\[
 \cZ(A)=\{X\in\Mat_n(\Q):AX=XA\}
\]
be the rational centralizer, and let
\[
 [A,\Mat_n(\Q)]=\{AY-YA:Y\in\Mat_n(\Q)\}
\]
be the image of the commutator map.  Under the rational semisimplicity hypotheses below, these
are complementary subspaces.

\begin{lemma}[Rational centralizer--commutator splitting]
\label{lem:centralizercommutator}
Suppose $A\in\Mat_n(\Q)$ is semisimple, all its eigenvalues
$\lambda_1,\dots,\lambda_s$ lie in $\Q$, and its spectral projectors $P_i$ lie in
$\Mat_n(\Q)$.  Then
\begin{equation}
 \Mat_n(\Q)=\cZ(A)\oplus[A,\Mat_n(\Q)].
 \label{eq:centralizerdecomp}
\end{equation}
The two projections are
\begin{align}
 H_{\mathrm c}&=\sum_iP_iHP_i,
 &
 H_{\mathrm o}&=\sum_{i\ne j}P_iHP_j.
 \label{eq:centraloff}
\end{align}
Moreover,
\begin{equation}
 H_{\mathrm o}=[A,Y_H],
 \qquad
 Y_H=\sum_{i\ne j}\frac1{\lambda_i-\lambda_j}P_iHP_j.
 \label{eq:YH}
\end{equation}
\end{lemma}

\begin{proof}
The block decomposition $H=\sum_{i,j}P_iHP_j$ is rational.  Its diagonal part commutes with
$A$.  For $i\ne j$,
\[
 [A,P_iHP_j]=(\lambda_i-\lambda_j)P_iHP_j,
\]
which proves \eqref{eq:YH}.  Finally, a commutator has zero diagonal spectral blocks, so the
intersection of the two summands is zero.
\end{proof}

\begin{lemma}[Commutator covariance]
\label{lem:commutatorcovariance}
For every holomorphic $f$ and every $Y$,
\begin{equation}
 Df_A[[A,Y]]=[f(A),Y].
 \label{eq:commutatorcovariance}
\end{equation}
\end{lemma}

\begin{proof}
Differentiate at $t=0$ the similarity identity
\[
 f(e^{-tY}Ae^{tY})=e^{-tY}f(A)e^{tY}.
\]
Equivalently, verify the identity termwise in \eqref{eq:divideddifference}.
\end{proof}

This covariance explains why commutator directions carry no new information: they are already
forced by ordinary similarity of the known value $f(A)$.

\subsection{The classification theorem}

\begin{theorem}[Exact first-order rationality locus]
\label{thm:rationalitylocus}
Let $A\in\Mat_n(\Q)$ be semisimple with positive rational eigenvalues
$\lambda_1,\dots,\lambda_s$ and rational spectral projectors.  Let $x\in\R\setminus\Q$ and
assume
\begin{equation}
 \lambda_i^x\in\Q\qquad(1\le i\le s).
 \label{eq:rationalspectralvalues}
\end{equation}
Then for every $H\in\Mat_n(\Q)$,
\begin{equation}
 Df_{x,A}[H]\in\Mat_n(\Q)
 \quad\Longleftrightarrow\quad
 H\in[A,\Mat_n(\Q)].
 \label{eq:rationalityiffcommutator}
\end{equation}
Equivalently,
\begin{equation}
 \cE_A(H)^x\in\Mat_{2n}(\Q)
 \quad\Longleftrightarrow\quad
 \cE_A(H)\text{ is rationally similar to }A\oplus A.
 \label{eq:blockrationalityiffsplit}
\end{equation}
\end{theorem}

\begin{proof}
Write $H=H_{\mathrm c}+H_{\mathrm o}$ as in
Lemma~\ref{lem:centralizercommutator}.  The off-diagonal part of
\eqref{eq:powerderivative} is rational because each coefficient
\[
 \frac{\lambda_i^x-\lambda_j^x}{\lambda_i-\lambda_j}
\]
is rational.  The centralizer part is
\begin{equation}
 Df_{x,A}[H_{\mathrm c}]
 =x\sum_i\lambda_i^{x-1}P_iHP_i
 =xC_H,
 \label{eq:xCH}
\end{equation}
where $C_H\in\Mat_n(\Q)$.  If $Df_{x,A}[H]$ is rational, then $xC_H$ is rational.
Since $x$ is irrational, every rational entry of $C_H$ must vanish, so $C_H=0$.  Multiplying
on the left and right by $P_i$ gives
\[
 \lambda_i^{x-1}P_iHP_i=0.
\]
The scalar is nonzero, hence every diagonal spectral block vanishes.  Thus
$H=H_{\mathrm o}\in[A,\Mat_n(\Q)]$.

Conversely, if $H=[A,Y]$, Lemma~\ref{lem:commutatorcovariance} gives
\[
 Df_{x,A}[H]=[A^x,Y]\in\Mat_n(\Q),
\]
because $A^x=\sum_i\lambda_i^xP_i$ is rational.

For the similarity statement, if $H=[A,Y]$, then
\begin{equation}
 \cE_A(H)=
 \begin{pmatrix}I&-Y\\0&I\end{pmatrix}
 \begin{pmatrix}A&0\\0&A\end{pmatrix}
 \begin{pmatrix}I&Y\\0&I\end{pmatrix}.
 \label{eq:splitsimilarity}
\end{equation}
Conversely, write $H=H_{\mathrm c}+[A,Y_H]$ by
Lemma~\ref{lem:centralizercommutator}.  The same block-unipotent calculation conjugates
$\cE_A(H)$ to $\cE_A(H_{\mathrm c})$.  On the $\lambda_i$-primary summand of the latter, the
nilpotent part is
$\left(\begin{smallmatrix}0&P_iHP_i\\0&0\end{smallmatrix}\right)$.  If $\cE_A(H)$ is similar
to the semisimple matrix $A\oplus A$, every such nilpotent part must vanish.  Hence
$H_{\mathrm c}=0$, so $H$ is a commutator.
\end{proof}

\begin{corollary}[A single non-split extension closes \AEname{}]
\label{cor:oneselfextension}
Assume \eqref{eq:baseline}.  Let $A$ satisfy the hypotheses of
Theorem~\ref{thm:rationalitylocus}, with all eigenvalues in $\Gamma$, and let
$H\in\Mat_n(\Q)$ have a nonzero centralizer component.  If
\[
 \cE_A(H)^x\in\Mat_{2n}(\Q),
\]
then $x\in\Z$.
\end{corollary}

\begin{proof}
Every eigenvalue $\lambda_i=2^{a_i}3^{b_i}$ has
$\lambda_i^x=M^{a_i}N^{b_i}\in\Q$.  If $x$ were irrational,
Theorem~\ref{thm:rationalitylocus} would force the centralizer component of $H$ to vanish.
Thus $x\in\Q$, and rational descent from $2^x\in\Q$ gives $x\in\Z$.
\end{proof}

\begin{remark}
The prior Jordan criterion corresponds to the special case $A=[b]$ and $H=[1]$.  Theorem
\ref{thm:rationalitylocus} is stronger in two ways: it treats arbitrary rational semisimple
bases and it proves necessity as well as sufficiency.  It identifies all false positives:
they are exactly the rationally split extensions.
\end{remark}

\section{The extension quotient and the missing scalar}
\label{sec:extensionquotient}

\subsection{Self-extension space}

Define the first-order rational self-extension quotient
\begin{equation}
 \Ext_A^{1,\mathrm{lin}}
 :=\Mat_n(\Q)/[A,\Mat_n(\Q)].
 \label{eq:extquotient}
\end{equation}
This notation is intentionally elementary: no derived-category machinery is needed.  Two block
extensions $\cE_A(H)$ and $\cE_A(H')$ are identified when a rational block-unipotent change of
basis changes one into the other, equivalently when $H-H'$ is a commutator.

If $A$ has spectral multiplicities $m_1,\dots,m_s$, then
\begin{equation}
 \dim_\Q \Ext_A^{1,\mathrm{lin}}=\sum_{i=1}^s m_i^2,
 \label{eq:extdimension}
\end{equation}
because the quotient is canonically represented by the block centralizer
$\bigoplus_i\End_\Q(P_i\Q^n)$.

\begin{theorem}[Power acts as $x$ times a rational automorphism]
\label{thm:quotientaction}
Under the hypotheses of Theorem~\ref{thm:rationalitylocus}, the derivative induces a map
\begin{equation}
 \overline{Df}_{x,A}:\Ext_A^{1,\mathrm{lin}}
 \longrightarrow \Ext_{A^x}^{1,\mathrm{lin}}.
 \label{eq:quotientmap}
\end{equation}
The centralizers and commutator spaces of $A$ and $A^x$ coincide, and after the natural
identification of the two quotients,
\begin{equation}
 \overline{Df}_{x,A}([H])
 =x\,[A^{x-1}H].
 \label{eq:xautomorphism}
\end{equation}
Multiplication by $A^{x-1}=A^xA^{-1}\in\Mat_n(\Q)$ is an invertible rational linear map on
the quotient.  Hence the only non-rational scalar in the induced action is $x$.
\end{theorem}

\begin{proof}
Lemma~\ref{lem:commutatorcovariance} shows that commutators map to commutators, so the map is
well-defined.  Since $x>0$ in the substantive case, the map $t\mapsto t^x$ is injective on the
positive spectrum; $A$ and $A^x$ therefore have the same spectral projectors and the same
block-diagonal/off-diagonal decomposition.  Modulo the off-diagonal commutator space,
\eqref{eq:powerderivative} becomes
\[
 [Df_{x,A}[H]]=x[A^{x-1}H].
\]
The matrix $A^{x-1}=A^xA^{-1}$ is rational and invertible.  Left multiplication by it preserves
each spectral block and the off-diagonal subspace, so it induces a rational automorphism of the
quotient.
\end{proof}

\begin{resultbox}
The original scalar problem can therefore be restated as follows.  The hypotheses give the
value-level map $A\mapsto A^x$ on every semisimple rational representation of the $3$-smooth
torus.  To prove $x\in\Z$, it is enough to show that this value-level map preserves one nonzero
rational first-order self-extension class.  On every such class, the induced scalar is $x$.
\end{resultbox}

\subsection{Basis-free extraction by traces and matrix coefficients}

The quotient theorem gives several concrete extraction formulas.  They are useful both for
paper arguments and for formalization because they avoid choosing eigenvectors.

\begin{proposition}[Weighted trace identity]
\label{prop:weightedtrace}
Under the hypotheses of Theorem~\ref{thm:quotientaction}, for every integer $m\ge0$ and every
$H\in\Mat_n(\Q)$,
\begin{equation}
 \tr\!\left(A^{m+1-x}Df_{x,A}[H]\right)
 =x\,\tr(A^mH).
 \label{eq:weightedtrace}
\end{equation}
Here
\[
 A^{m+1-x}=A^{m+1}(A^x)^{-1}\in\Mat_n(\Q).
\]
\end{proposition}

\begin{proof}
Only diagonal spectral blocks contribute to a trace.  On the $i$-th diagonal block,
$Df_{x,A}$ acts by $x\lambda_i^{x-1}$, while $A^{m+1-x}$ acts by
$\lambda_i^{m+1-x}$.  Their product is $x\lambda_i^m$.  Summing the block traces gives
\eqref{eq:weightedtrace}.
\end{proof}

\begin{corollary}[Moment extraction]
\label{cor:momentextraction}
If $Df_{x,A}[H]$ is rational and $\tr(A^mH)\ne0$ for some $m\ge0$, then $x\in\Q$.
Under \eqref{eq:baseline}, this implies $x\in\Z$.
\end{corollary}

\begin{proof}
The left side of \eqref{eq:weightedtrace} is rational and the right side is a nonzero rational
multiple of $x$.
\end{proof}

\begin{proposition}[Finite detection for simple spectrum]
\label{prop:finitedetection}
Suppose in addition that $A$ has $n$ distinct rational eigenvalues.  If the diagonal spectral
component of $H$ is nonzero, then
\[
 \tr(A^mH)\ne0
\]
for at least one $m\in\{0,1,\dots,n-1\}$.
\end{proposition}

\begin{proof}
Choose a rational eigenbasis.  Let $h_i$ be the $i$-th diagonal entry of $H$.  The equations
$\sum_i\lambda_i^m h_i=0$ for $0\le m<n$ have Vandermonde coefficient matrix and therefore
force every $h_i=0$.
\end{proof}

\begin{warning}
With spectral multiplicities, traces detect only the traces of the centralizer blocks; a
nonzero traceless block may escape all scalar moments.  The full theorem is nevertheless
recovered by rational matrix coefficients $u^TP_iDf_{x,A}[H]P_iv$, or simply by the block
projection itself.  No trace-only claim is made in that generality.
\end{warning}

\section{The complete \texorpdfstring{$2$--$3$}{2--3} block model}
\label{sec:twothreemodel}

Take
\begin{equation}
 A=\begin{pmatrix}2&0\\0&3\end{pmatrix},
 \qquad
 H=\begin{pmatrix}a&b\\c&d\end{pmatrix}\in\Mat_2(\Q).
 \label{eq:A23H}
\end{equation}
Then $A^x=\diag(M,N)$, and Theorem~\ref{thm:semisimplederivative} gives the exact formula
\begin{equation}
 Df_{x,A}[H]=
 \begin{pmatrix}
 \dfrac{xM}{2}a&(N-M)b\\[0.8ex]
 (N-M)c&\dfrac{xN}{3}d
 \end{pmatrix}.
 \label{eq:explicit23derivative}
\end{equation}
Consequently
\begin{equation}
 \begin{pmatrix}
 2&0&a&b\\
 0&3&c&d\\
 0&0&2&0\\
 0&0&0&3
 \end{pmatrix}^{\!x}
 =
 \begin{pmatrix}
 M&0&\frac{xM}{2}a&(N-M)b\\
 0&N&(N-M)c&\frac{xN}{3}d\\
 0&0&M&0\\
 0&0&0&N
 \end{pmatrix}.
 \label{eq:explicit4by4}
\end{equation}

\begin{theorem}[Explicit rationality test]
\label{thm:explicit23test}
Assume \eqref{eq:baseline} and $x\notin\Q$.  For $H$ as in \eqref{eq:A23H},
\begin{equation}
 \cE_A(H)^x\in\Mat_4(\Q)
 \quad\Longleftrightarrow\quad
 a=d=0.
 \label{eq:aanddzero}
\end{equation}
The matrices with $a=d=0$ are precisely the commutators $[A,Y]$.
\end{theorem}

\begin{proof}
The off-diagonal entries in \eqref{eq:explicit23derivative} are rational.  Since $M,N$ are
nonzero rational numbers and $x$ is irrational, the diagonal entries are rational precisely
when $a=d=0$.  Directly,
\[
 [A,\begin{pmatrix}u&v\\w&z\end{pmatrix}]
 =\begin{pmatrix}0&-v\\w&0\end{pmatrix},
\]
which is exactly the off-diagonal subspace.
\end{proof}

\begin{corollary}[Minimal concrete closing target]
\label{cor:minimalclosing}
It would suffice to derive from the original scalar hypotheses the rationality of the $x$-th
power of any one of the following matrices:
\[
 \begin{pmatrix}
 2&0&1&0\\0&3&0&0\\0&0&2&0\\0&0&0&3
 \end{pmatrix},
 \qquad
 \begin{pmatrix}
 2&0&0&0\\0&3&0&1\\0&0&2&0\\0&0&0&3
 \end{pmatrix}.
\]
The exposed entry is respectively $xM/2$ or $xN/3$, and rational descent then gives
$x\in\Z$.
\end{corollary}

\subsection{False positives are exactly split directions}

For any $u,v\in\Q$, the off-diagonal direction
$H=\left(\begin{smallmatrix}0&u\\v&0\end{smallmatrix}\right)$ produces a rational powered
block.  This is not partial evidence for a derivative.  Equation~\eqref{eq:splitsimilarity}
shows that it is just a change of rational basis of $A\oplus A$.  The classification theorem
is therefore also a debugging tool: any proposed finite-dimensional construction should first
be projected to the quotient \eqref{eq:extquotient}.  If its class vanishes, no amount of
entrywise complexity can reveal $x$.

\section{Why toric functors create multiplicity but not confluence}
\label{sec:toricnogo}

\subsection{Weight-space formulation}

Let $T=\mathbb G_m^2$ over $\Q$.  A finite-dimensional rational representation of $T$ splits
as a direct sum of weight spaces
\begin{equation}
 V=\bigoplus_{(a,b)\in\Z^2}V_{a,b},
 \qquad
 (s,t)\cdot v=s^at^b v
 \quad(v\in V_{a,b}).
 \label{eq:weightsplit}
\end{equation}
This is not merely diagonalizability at one matrix: a split torus is linearly reductive in
characteristic zero, so every rational subrepresentation and every quotient again decomposes
into weights.

Evaluating at $(2,3)$ gives eigenvalue $2^a3^b$ on $V_{a,b}$; evaluating at $(M,N)$ gives
$M^aN^b=(2^a3^b)^x$.  Unique factorization implies
\begin{equation}
 2^a3^b=2^{a'}3^{b'}
 \quad\Longleftrightarrow\quad
 (a,b)=(a',b').
 \label{eq:weightinjective}
\end{equation}
The same remains true after applying $x>0$.

\begin{theorem}[Toric-functor no-go]
\label{thm:toricnogo}
Assume $x>0$.  Start with finite-dimensional rational $T$-representations and apply any finite combination of
\begin{enumerate}
\item direct sums and tensor products;
\item duals;
\item symmetric, exterior, or more general Schur functors;
\item rational $T$-stable subspaces and rational quotients.
\end{enumerate}
After evaluation at $(2,3)$ or at $(M,N)$, the resulting operator is semisimple.  Repeated
eigenvalues can arise only from repeated occurrences of the same character; the operator on
each repeated eigenspace is scalar.  In particular, every self-extension arising within the
toric representation category is split and maps to the zero class in
\eqref{eq:extquotient}.
\end{theorem}

\begin{proof}
Every listed operation stays inside the semisimple tensor category of rational representations
of the split torus.  Thus the output is a direct sum of character spaces.  Evaluation sends a
character $(a,b)$ to the positive scalar $2^a3^b$ or $M^aN^b$.  By
\eqref{eq:weightinjective}, distinct characters do not collide under either evaluation.
Multiplicity can therefore only mean that one character occurs several times, and the torus
element acts as a scalar on that entire multiplicity space.  In particular, every
self-extension constructed inside this toric representation category splits; its image in
$\Ext_A^{1,\mathrm{lin}}$ is the zero class.  (The ambient quotient
$\Ext_A^{1,\mathrm{lin}}$ itself is generally nonzero.)
\end{proof}

\begin{corollary}[Spectral collision is not confluence]
\label{cor:collisionnotconfluence}
Tensor-product identities such as
\[
 (2)(3)=(3)(2)
\]
produce a repeated eigenvalue $6$, but the two copies are copies of the same weight and remain
a direct sum.  They do not produce the derivative coefficient $x6^{x-1}$.
\end{corollary}

\begin{remark}
This theorem does not say that tensor constructions are useless for every purpose.  They can
amplify valuations, create determinant identities, and expose multiplicities.  It says
something narrower and exact: without a non-semisimple operation or a singular limiting
process, they cannot cross from value data to first-jet data.
\end{remark}

\subsection{Where a successful construction must leave the torus category}

There are three mathematically distinct exits:

\begin{description}
\item[An actual extension.] Produce a rational matrix with a nonzero class in
$\Ext_A^{1,\mathrm{lin}}$ and prove its $x$-th power rational.

\item[A singular degeneration.] Approach a repeated spectral point through distinct rational
spectral points while controlling denominators strongly enough that the limiting jet is
arithmetic.

\item[An affine relation.] Add translation or another non-toric operation.  The relation
$3=2+1$ is the smallest available additive input.
\end{description}

The remainder of the report studies the last two exits.

\section{Integral lattice gluing without a jet}
\label{sec:newtonchain}

\subsection{A \texorpdfstring{$3\times3$}{3 x 3} Newton chain at \texorpdfstring{$2,3,4$}{2, 3, 4}}

Consider the integral upper-bidiagonal matrix
\begin{equation}
 T=
 \begin{pmatrix}
 2&1&0\\
 0&3&1\\
 0&0&4
 \end{pmatrix}.
 \label{eq:T234}
\end{equation}
Its eigenvalues are distinct, so it is semisimple over $\Q$, but its integral lattice is
nontrivially glued across the three eigenspaces.

\begin{theorem}[Integral Newton-chain power]
\label{thm:T234power}
Under \eqref{eq:baseline},
\begin{equation}
 T^x=
 \begin{pmatrix}
 M&N-M&\dfrac{M^2-2N+M}{2}\\[1ex]
 0&N&M^2-N\\
 0&0&M^2
 \end{pmatrix}
 \in\Mat_3(\Z).
 \label{eq:T234power}
\end{equation}
\end{theorem}

\begin{proof}
For an upper-bidiagonal matrix with distinct diagonal entries, functional calculus gives first
divided differences on the first superdiagonal and the second divided difference on the
second superdiagonal.  Here
\begin{align*}
 f_x[2,3]&=N-M,\\
 f_x[3,4]&=M^2-N,\\
 f_x[2,3,4]
 &=\frac{f_x[3,4]-f_x[2,3]}{4-2}
 =\frac{M^2-2N+M}{2}.
\end{align*}
The final numerator is even because $M^2+M$ is even and $2N$ is even.  Thus every entry is an
integer.
\end{proof}

\begin{proposition}[The lattice gluing is genuinely integral]
\label{prop:index2}
The matrix $T$ is diagonalizable over $\Q$ but not conjugate to $\diag(2,3,4)$ by an element of
$\operatorname{GL}_3(\Z)$.  The lattice generated by primitive integral eigenvectors has index
$2$ in $\Z^3$.
\end{proposition}

\begin{proof}
Primitive eigenvectors for $2,3,4$ are respectively
\[
 v_2=(1,0,0)^T,\qquad
 v_3=(1,1,0)^T,\qquad
 v_4=(1,2,2)^T.
\]
Their column determinant is $2$.  Any integral eigenbasis has the same primitive lines and
therefore determinant divisible by $2$, so no unimodular eigenbasis exists.
\end{proof}

This example matters because it rules out an overly optimistic inference: ``if a rational
semisimple matrix is integrally glued, then integrality of its irrational power should force
an integer exponent.''  The original hypotheses themselves make \eqref{eq:T234power}
integral.  Integral gluing is not yet a self-extension.

\subsection{A concrete repeated-spectrum false positive}

The tensor square $T\otimes T$ is integral and semisimple.  Its eigenvalue multiset is
\[
 4,6,8,6,9,12,8,12,16,
\]
so $6,8,12$ occur with multiplicity.  Moreover,
\begin{equation}
 (T\otimes T)^x=T^x\otimes T^x\in\Mat_9(\Z).
 \label{eq:tensorpower}
\end{equation}
The equality follows after rational diagonalization, since the eigenvalues of a tensor product
are products and $(\lambda\mu)^x=\lambda^x\mu^x$ for positive scalars.

\begin{audit}[Collision without confluence]
The matrix $T\otimes T$ has repeated eigenvalues, a nontrivial integral lattice, and an integral
$x$-th power.  Nevertheless it is semisimple, and every repeated eigenspace carries a scalar
operator.  Its first-order extension class is zero.  Thus none of the three visible features
--- repeated spectrum, integral gluing, or integral powered entries --- substitutes for a
Jordan jet.
\end{audit}

\section{Split matrices can degenerate to a Jordan block}
\label{sec:splitdegeneration}

The previous no-go theorem is exact at every finite stage, but semisimplicity is not closed
under singularly changing bases.  This produces a useful model of the analytic limit and a
warning about arithmetic compactness.  Such integral collisions exist: if
$p_j/q_j$ are continued-fraction convergents to $\log3/\log2$, then one may take
$a_j=2^{p_j}$ and $b_j=3^{q_j}$ along either side of the approximation.

\begin{theorem}[Unbounded integral conjugators recover the first jet]
\label{thm:splitdegeneration}
Let $a_j,b_j$ be positive elements of the integral $3$-smooth semigroup
$\{2^u3^v:u,v\in\Nzero\}$ such that
\begin{equation}
 r_j:=\frac{b_j}{a_j}\longrightarrow1,
 \qquad b_j\ne a_j.
 \label{eq:relativecollision}
\end{equation}
Choose integers $k_j$ satisfying
\begin{equation}
 k_j(r_j-1)\longrightarrow1.
 \label{eq:kchoice}
\end{equation}
For example, take an integer nearest to $(r_j-1)^{-1}$.  Define
\begin{equation}
 S_j=
 \begin{pmatrix}1&-k_j\\0&1\end{pmatrix},
 \qquad
 B_j=S_j^{-1}
 \begin{pmatrix}a_j&0\\0&b_j\end{pmatrix}S_j.
 \label{eq:Bj}
\end{equation}
Then $B_j\in\Mat_2(\Z)$, $B_j^x\in\Mat_2(\Z)$, and
\begin{align}
 a_j^{-1}B_j&\longrightarrow
 J:=\begin{pmatrix}1&1\\0&1\end{pmatrix},
 \label{eq:Bjlimit}\\
 a_j^{-x}B_j^x&\longrightarrow
 J^x=\begin{pmatrix}1&x\\0&1\end{pmatrix}.
 \label{eq:Bjxlimit}
\end{align}
\end{theorem}

\begin{proof}
A direct multiplication gives
\begin{equation}
 B_j=
 \begin{pmatrix}
 a_j&k_j(b_j-a_j)\\0&b_j
 \end{pmatrix}.
 \label{eq:Bjexplicit}
\end{equation}
It is integral and integrally similar to a diagonal matrix.  Since $a_j^x$ and $b_j^x$ are
integers,
\begin{equation}
 B_j^x=
 \begin{pmatrix}
 a_j^x&k_j(b_j^x-a_j^x)\\0&b_j^x
 \end{pmatrix}
 \in\Mat_2(\Z).
 \label{eq:Bjxexplicit}
\end{equation}
After division by $a_j$, the upper-right entry of \eqref{eq:Bjexplicit} is
$k_j(r_j-1)$, which tends to $1$ by \eqref{eq:kchoice}.  For the power, the normalized
upper-right entry is
\[
 k_j(r_j^x-1)
 =k_j(r_j-1)\frac{r_j^x-1}{r_j-1}
 \longrightarrow x.
\]
The diagonal entries converge to $1$, proving both limits.
\end{proof}

\begin{remark}
The theorem is an exact matrix avatar of transported secants.  At every finite stage the
extension is split by the integral matrix $S_j$, but $\|S_j\|\to\infty$.  The limit is a
non-split Jordan block, and its powered limit contains $x$.  The obstruction is therefore not
analytic existence of the jet; it is the lack of a bounded arithmetic lattice in which the
limit can be taken.
\end{remark}

\begin{corollary}[Symmetric powers expose the full limiting jet tower]
\label{cor:sympowerjet}
For each fixed $d\ge1$, apply the $d$-th symmetric power to
Theorem~\ref{thm:splitdegeneration}.  Then
\[
 a_j^{-dx}\bigl(\Sym^d(B_j)\bigr)^x
 =a_j^{-dx}\Sym^d(B_j^x)
 =\Sym^d\!\left(a_j^{-x}B_j^x\right)
 \longrightarrow
 \Sym^d\!\begin{pmatrix}1&x\\0&1\end{pmatrix}.
\]
In the monomial basis, the limiting entries are integer binomial multiples of
$1,x,x^2,\dots,x^d$.
\end{corollary}

\begin{proof}
The identity holds after diagonalizing $B_j$: both sides have the same eigenvectors and
weights $(a_j^{d-k}b_j^k)^x$.  Symmetric power is polynomial in the matrix entries, so it
commutes with the displayed limits.
\end{proof}

\begin{warning}
Corollary~\ref{cor:sympowerjet} does not contradict the toric no-go theorem.  Every finite
$j$ remains semisimple and split.  The higher jets appear only after an unbounded change of
basis and non-integral normalizations by powers of $a_j^x$.  The denominators escaping through
those normalizations are exactly the arithmetic difficulty.
\end{warning}

\section{Affine collision at the additive relation \texorpdfstring{$3=2+1$}{3 = 2 + 1}}
\label{sec:affinecollision}

The torus category knows multiplication but not addition.  The smallest way to leave it while
remaining close to the original bases is to use the additive relation $3=2+1$.

\subsection{A rational secant family converging to the decisive Jordan block}

For $q\in\Q_{>0}$ define
\begin{equation}
 C(q)=\begin{pmatrix}q&1\\0&3\end{pmatrix}.
 \label{eq:Cq}
\end{equation}
For $q\ne3$, the eigenvalues are distinct and
\begin{equation}
 C(q)^x=
 \begin{pmatrix}
 q^x&\dfrac{q^x-N}{q-3}\\[1.2ex]
 0&N
 \end{pmatrix}.
 \label{eq:Cqpower}
\end{equation}
At $q=3$, the removable divided difference becomes a derivative:
\begin{equation}
 C(3)^x=
 \begin{pmatrix}
 N&x3^{x-1}\\0&N
 \end{pmatrix}.
 \label{eq:C3power}
\end{equation}

If $q\in\Gamma$, then $q^x=\rho_x(q)\in\Q$, so every entry of
\eqref{eq:Cqpower} is rational.  Since $\Gamma$ is dense in logarithmic scale, there are
$q_j\in\Gamma\setminus\{3\}$ with $q_j\to3$.  Thus the decisive irrationality detector
$x3^{x-1}$ is the real limit of rational matrix entries supplied entirely by the original
value data.

\begin{proposition}[Affine collision criterion]
\label{prop:affinecriterion}
Assume \eqref{eq:baseline}.  Any one of the following statements implies $x\in\Z$:
\begin{enumerate}
\item $C(3)^x\in\Mat_2(\Q)$;
\item the rational secants $(q^x-N)/(q-3)$, for $q\in\Gamma$ approaching $3$, are eventually
constant;
\item these secants lie in a fixed discrete subgroup $D^{-1}\Z\subset\Q$ in a neighborhood of
$3$;
\item their reduced denominators remain bounded along some sequence $q_j\to3$.
\end{enumerate}
\end{proposition}

\begin{proof}
The first statement makes $x3^{x-1}$ rational, hence
$x=3(x3^{x-1})/N\in\Q$.  For the others, the secants converge in $\R$ to
$x3^{x-1}$.  An eventually constant rational sequence has a rational limit.  A convergent
sequence in a fixed discrete subgroup is eventually constant.  Reduced denominators bounded
by $D$ put all terms in $(\lcm(1,\dots,D))^{-1}\Z$, reducing to the discrete-subgroup case.
Rational descent finishes.
\end{proof}

\begin{target}[Arithmetic compactness at an affine collision]
\label{target:affinecompactness}
Prove any denominator bound, congruence stabilization, adelic compactness principle, or exact
lattice realization strong enough to trigger Proposition~\ref{prop:affinecriterion} for
$q\in\Gamma$ approaching $3$.
\end{target}

This target is genuinely non-toric.  The parameter $q$ moves through a multiplicative group,
but the collision point is selected by the affine difference $q-3$.  The denominator
$q-3$ is where additive arithmetic enters.

\subsection{Integral orders and the conductor defect}

The secant denominator has an exact lattice interpretation.  Let $a,b,A,B\in\Z$ with
$a\ne b$.  For an integer $d\ge1$, define the order
\begin{equation}
 \cO_d=\{(u,v)\in\Z^2:u\equiv v\pmod d\}\subset\Q\times\Q.
 \label{eq:Od}
\end{equation}
Multiplication by $(a,b)$ preserves $\cO_d$ exactly when $d\mid a-b$.  In the basis
$(1,1),(0,d)$ it has matrix
\begin{equation}
 \begin{pmatrix}
 a&0\\(b-a)/d&b
 \end{pmatrix}.
 \label{eq:orderaction}
\end{equation}
Multiplication by $(A,B)$ preserves the same order exactly when $d\mid A-B$.

\begin{definition}[Common conductor and conductor defect]
For a pair of integral $3$-smooth numbers $a,b$, put
\begin{equation}
 a_x=a^x\in\Z,
 \qquad b_x=b^x\in\Z,
 \qquad
 g_x(a,b)=\gcd(a-b,a_x-b_x),
 \label{eq:commonconductor}
\end{equation}
and define
\begin{equation}
 \delta_x(a,b)=\frac{|a-b|}{g_x(a,b)}.
 \label{eq:conductordefect}
\end{equation}
\end{definition}

The integer $g_x(a,b)$ is the largest conductor $d$ for which both the input element $(a,b)$
and its powered output $(a_x,b_x)$ preserve the same glued order $\cO_d$.  The defect
$\delta_x(a,b)$ is exactly the reduced denominator of the ordinary divided difference
$(a_x-b_x)/(a-b)$.

\begin{proposition}[Order-preservation reformulation]
\label{prop:orderreformulation}
For integral $3$-smooth $a\ne b$, the following are equivalent:
\begin{enumerate}
\item $(a^x-b^x)/(a-b)\in\Z$;
\item $a-b\mid a^x-b^x$;
\item the maximal input order $\cO_{|a-b|}$ is also preserved by $(a^x,b^x)$;
\item the matrix $\left(\begin{smallmatrix}a&1\\0&b\end{smallmatrix}\right)^x$ is integral.
\end{enumerate}
More generally, the reduced denominator in (1) is $\delta_x(a,b)$.
\end{proposition}

\begin{proof}
The equivalence of the first three statements follows immediately from the order calculation.
The fourth follows from the two-point divided-difference formula.  Reducing the fraction
$(a^x-b^x)/(a-b)$ gives denominator $|a-b|/\gcd(a-b,a^x-b^x)$.
\end{proof}

\begin{remark}
This formulation suggests a concrete computational reconnaissance program: search close
$3$-smooth pairs not merely for a small real gap, but for anomalously small conductor defect.
The real approximation and the lattice preservation must occur simultaneously.  A small gap
without common conductor reproduces the old denominator wall; a large common conductor is the
arithmetic signature of genuine confluence.
\end{remark}

\section{Symmetric transported secants}
\label{sec:symmetricsecants}

The transported secants in the unified report use
\[
 \frac{r^x-1}{r-1}
 =\frac{e^{xh}-1}{e^h-1},
 \qquad h=\log r,
\]
and approximate $x$ with a linear error in $h$.  A symmetric quotient cancels every odd error
term automatically.

\begin{definition}[Symmetric secant]
For $r\in\Gamma\setminus\{1\}$ define
\begin{equation}
 J_x(r)=\frac{r^x-r^{-x}}{r-r^{-1}}.
 \label{eq:Jdefinition}
\end{equation}
Writing $r=e^h$ gives
\begin{equation}
 J_x(e^h)=\frac{\sinh(xh)}{\sinh h}.
 \label{eq:Jsinh}
\end{equation}
\end{definition}

Under \eqref{eq:baseline}, $r^x\in\Q$ and hence $J_x(r)\in\Q$.

\begin{theorem}[Second-order rational confluence]
\label{thm:secondorder}
For fixed real $x$,
\begin{equation}
 J_x(e^h)
 =x+\frac{x(x^2-1)}6h^2
  +\frac{x(3x^4-10x^2+7)}{360}h^4
  +O_x(h^6)
 \qquad(h\to0).
 \label{eq:Jexpansion}
\end{equation}
In particular,
\begin{equation}
 J_x(e^h)=x+O_x(h^2).
 \label{eq:Jsecondorder}
\end{equation}
\end{theorem}

\begin{proof}
Divide the Taylor series
\[
 \sinh(xh)=xh+\frac{x^3h^3}{6}+\frac{x^5h^5}{120}+O_x(h^7)
\]
by
\[
 \sinh h=h+\frac{h^3}{6}+\frac{h^5}{120}+O(h^7).
\]
Formal series division gives \eqref{eq:Jexpansion}.  Evenness in $h$ also shows directly that
all odd powers vanish.
\end{proof}

\begin{remark}
The coefficient of $h^4$ is included because it is a useful audit check.  For $x=2$,
$J_2(e^h)=2\cosh h=2+h^2+h^4/12+\cdots$, and the formula gives $1/12$.
For $x=1$, every correction vanishes and $J_1\equiv1$.
\end{remark}

\subsection{Exact arithmetic formula}

Let $r=B/A$ and $r^x=D/C$, where $A,B,C,D$ are positive integers.  No coprimality assumption
is needed for the following unreduced identity:
\begin{equation}
 J_x(r^k)
 =\frac{A^kB^k(D^{2k}-C^{2k})}
 {C^kD^k(B^{2k}-A^{2k})}.
 \label{eq:Jintegerformula}
\end{equation}
For the canonical near-collisions $r=3^q/2^p$, one may take
\[
 A=2^p,
 \quad B=3^q,
 \quad C=M^p,
 \quad D=N^q.
\]
Formula~\eqref{eq:Jintegerformula} makes the height problem visible: even when
$|\log r|$ is of order $1/q$, the unreduced numerator and denominator have exponential size in
$p+q$.

\section{Richardson extrapolation to arbitrary even order}
\label{sec:richardson}

\subsection{The coefficients}

For integers $m\ge1$ and $1\le k\le m$, define
\begin{equation}
 c_{m,k}=(-1)^{k-1}\frac{2(m!)^2}{(m-k)!(m+k)!}.
 \label{eq:cmk}
\end{equation}
These are the Lagrange coefficients for evaluating at $0$ a polynomial given at the nodes
$1^2,2^2,\dots,m^2$.

\begin{lemma}[Moment cancellation]
\label{lem:momentcancellation}
For every $m\ge1$,
\begin{align}
 \sum_{k=1}^m c_{m,k}&=1,
 \label{eq:moment0}\\
 \sum_{k=1}^m c_{m,k}k^{2j}&=0
 \qquad(1\le j<m).
 \label{eq:momentvanish}
\end{align}
\end{lemma}

\begin{proof}
Let $L_k(t)$ be the Lagrange basis polynomial for the nodes $1^2,\dots,m^2$.  Its value at zero
is
\begin{align*}
 L_k(0)
 &=\prod_{\substack{1\le\ell\le m\\\ell\ne k}}
   \frac{-\ell^2}{k^2-\ell^2}\\
 &=(-1)^{k-1}\frac{2(m!)^2}{(m-k)!(m+k)!}
 =c_{m,k}.
\end{align*}
Applying the interpolation identity $P(0)=\sum_kL_k(0)P(k^2)$ to $P(t)=1$ gives
\eqref{eq:moment0}; applying it to $P(t)=t^j$, $1\le j<m$, gives
\eqref{eq:momentvanish}.
\end{proof}

The first rows are
\begin{align*}
 (c_{1,k})&=(1),\\
 (c_{2,k})&=\left(\frac43,-\frac13\right),\\
 (c_{3,k})&=\left(\frac32,-\frac35,\frac1{10}\right),\\
 (c_{4,k})&=\left(\frac85,-\frac45,\frac8{35},-\frac1{35}\right).
\end{align*}

\subsection{The extrapolants}

\begin{definition}[Order-$2m$ rational jet approximant]
For $r\in\Gamma\setminus\{1\}$, define
\begin{equation}
 R_{m,x}(r)=\sum_{k=1}^m c_{m,k}J_x(r^k).
 \label{eq:Rdefinition}
\end{equation}
\end{definition}

Every term is rational under \eqref{eq:baseline}, so
\begin{equation}
 R_{m,x}(r)\in\Q.
 \label{eq:Rrational}
\end{equation}
Explicitly,
\begin{align}
 R_{1,x}(r)&=J_x(r),
 \label{eq:R1}\\
 R_{2,x}(r)&=\frac{4J_x(r)-J_x(r^2)}3,
 \label{eq:R2}\\
 R_{3,x}(r)&=\frac{15J_x(r)-6J_x(r^2)+J_x(r^3)}{10},
 \label{eq:R3}\\
 R_{4,x}(r)&=\frac{56J_x(r)-28J_x(r^2)+8J_x(r^3)-J_x(r^4)}{35}.
 \label{eq:R4}
\end{align}

\begin{theorem}[Arbitrary fixed even-order confluence]
\label{thm:richardsonorder}
For every fixed $x\in\R$ and $m\ge1$,
\begin{equation}
 R_{m,x}(e^h)=x+O_{x,m}(h^{2m})
 \qquad(h\to0).
 \label{eq:Richardsonerror}
\end{equation}
More precisely, if
\begin{equation}
 \frac{\sinh(xz)}{\sinh z}=x+\sum_{j\ge1}a_j(x)z^{2j}
 \label{eq:evenseries}
\end{equation}
near $z=0$, then
\begin{equation}
 R_{m,x}(e^h)-x
 =a_m(x)h^{2m}\sum_{k=1}^m c_{m,k}k^{2m}+O_{x,m}(h^{2m+2}).
 \label{eq:leadingRichardson}
\end{equation}
\end{theorem}

\begin{proof}
Insert \eqref{eq:evenseries} with $z=kh$ into \eqref{eq:Rdefinition}.  The constant term is
$x\sum_kc_{m,k}=x$.  For $1\le j<m$, the coefficient of $h^{2j}$ is
$a_j(x)\sum_kc_{m,k}k^{2j}=0$ by Lemma~\ref{lem:momentcancellation}.  The first term not
forced to vanish is $j=m$, which gives \eqref{eq:leadingRichardson} and hence
\eqref{eq:Richardsonerror}.
\end{proof}

\begin{proposition}[First surviving moment]
\label{prop:firstsurvivingmoment}
The first uncancelled moment is
\begin{equation}
 \sum_{k=1}^m c_{m,k}k^{2m}=(-1)^{m-1}(m!)^2.
 \label{eq:firstsurvivingmoment}
\end{equation}
\end{proposition}

\begin{proof}
Apply the degree-$m$ Lagrange interpolation remainder to $P(t)=t^m$ at $t=0$.  Since the
interpolating polynomial through the nodes $1^2,\dots,m^2$ has degree at most $m-1$,
\[
 P(0)-\sum_kc_{m,k}P(k^2)
 =\prod_{k=1}^m(0-k^2)=(-1)^m(m!)^2.
\]
As $P(0)=0$, rearrangement gives \eqref{eq:firstsurvivingmoment}.
\end{proof}

Combining Theorem~\ref{thm:richardsonorder} and
Proposition~\ref{prop:firstsurvivingmoment} gives the sharpened leading term
\begin{equation}
 R_{m,x}(e^h)-x
 =(-1)^{m-1}(m!)^2a_m(x)h^{2m}+O_{x,m}(h^{2m+2}).
 \label{eq:sharpleading}
\end{equation}

\subsection{Existence of rational approximants}

By continued fractions for $\log3/\log2$, there are infinitely many coprime positive pairs
$(p,q)$ with
\begin{equation}
 h_{p,q}=q\log3-p\log2\longrightarrow0.
 \label{eq:hpq}
\end{equation}
Taking
\[
 r_{p,q}=\frac{3^q}{2^p}\in\Gamma
\]
and applying Theorem~\ref{thm:richardsonorder} yields the following.

\begin{corollary}[Rational approximations of every fixed even order]
\label{cor:rationalapproximants}
Under \eqref{eq:baseline}, for every fixed $m\ge1$ there are explicitly computable rational
numbers
\begin{equation}
 R_{m,x}\!\left(\frac{3^q}{2^p}\right)
 \label{eq:explicitrationalapprox}
\end{equation}
with
\begin{equation}
 \left|R_{m,x}\!\left(\frac{3^q}{2^p}\right)-x\right|
 \ll_{x,m}|q\log3-p\log2|^{2m}.
 \label{eq:approxbound}
\end{equation}
For convergents $p/q$, the right side is $O_{x,m}(q^{-2m})$.
\end{corollary}

\begin{proof}
Rationality follows from \eqref{eq:Rrational}.  The analytic estimate is
Theorem~\ref{thm:richardsonorder}.  Continued-fraction convergents satisfy
$|q\log3-p\log2|\ll q^{-1}$.
\end{proof}

This is the strongest unconditional approximation statement obtained in this continuation.
It gives arbitrary fixed contact order at the missing jet while using only the rational values
forced by the original hypotheses.

\section{Denominator and height audit}
\label{sec:denaudit}

High-order contact is not by itself an irrationality proof.  This section gives the exact
reason, in a form suitable for deciding what kind of new lemma would be decisive.

\subsection{Naive height of the rational extrapolants}

Let $(p,q)$ be a convergent pair and use
\[
 A=2^p,
 \quad B=3^q,
 \quad C=M^p,
 \quad D=N^q.
\]
For fixed $m$, Formula~\eqref{eq:Jintegerformula} and the rational coefficients
\eqref{eq:cmk} give a completely explicit common denominator for
$R_{m,x}(B/A)$.  Since $p\asymp q$ along convergents, one obtains
\begin{equation}
 \log \den\!\left(R_{m,x}(B/A)\right)
 \le C_{x,m}q
 \label{eq:fixedmheight}
\end{equation}
for an effective constant $C_{x,m}$, before any favorable cancellation is used.  If $m$ is
allowed to vary, the crude product denominator gives
\begin{equation}
 \log \den\!\left(R_{m,x}(B/A)\right)
 =O_x(m^2q+m^2\log(m+1)).
 \label{eq:variablemheight}
\end{equation}
These are upper bounds, not claims that the reduced denominator is always this large.  The
possible cancellation is precisely the unexplored arithmetic resource.

For fixed $m$, Corollary~\ref{cor:rationalapproximants} supplies an error $O(q^{-2m})$, while
\eqref{eq:fixedmheight} allows denominators of size $\exp(Cq)$.  The reciprocal denominator
scale is exponentially small, far below the polynomial analytic error.  Consequently neither
ordinary rational spacing nor an effective irrationality measure can yet lock the sequence.

\subsection{Adjacent rational locking}

\begin{lemma}[Unit locking for two rational approximants]
\label{lem:unitlocking}
Let $P_j/Q_j$ and $P_{j+1}/Q_{j+1}$ be reduced rationals with positive denominators.  If
\begin{equation}
 \left|x-\frac{P_j}{Q_j}\right|
 +\left|x-\frac{P_{j+1}}{Q_{j+1}}\right|
 <\frac1{Q_jQ_{j+1}},
 \label{eq:unitlockingcondition}
\end{equation}
then
\[
 \frac{P_j}{Q_j}=\frac{P_{j+1}}{Q_{j+1}}.
\]
\end{lemma}

\begin{proof}
If the rationals were distinct, their distance would be at least $1/(Q_jQ_{j+1})$.  The
triangle inequality and \eqref{eq:unitlockingcondition} contradict this.
\end{proof}

\begin{criterion}[Richardson adjacent locking]
\label{crit:richardsonlocking}
Choose $r_j=e^{h_j}\in\Gamma$ with $h_j\to0$, orders $m_j\ge1$, and write the reduced rational
extrapolants as
\[
 R_{m_j,x}(r_j)=\frac{P_j}{Q_j}.
\]
Suppose explicit remainder bounds $E_j$ are available with
\[
 \left|R_{m_j,x}(r_j)-x\right|\le E_j
\]
and
\begin{equation}
 E_j+E_{j+1}<\frac1{Q_jQ_{j+1}}
 \label{eq:lockingtarget}
\end{equation}
for every sufficiently large $j$.  Then $x\in\Q$, hence $x\in\Z$ under
\eqref{eq:baseline}.
\end{criterion}

\begin{proof}
Lemma~\ref{lem:unitlocking} makes the rational sequence eventually constant.  Its real limit is
$x$ by Theorem~\ref{thm:richardsonorder}, so $x$ is that rational constant.
\end{proof}

The criterion is elementary, exact, and stronger than any appeal to a generic irrationality
measure.  It identifies the desired denominator threshold without hiding constants.

\subsection{Irrationality-measure locking}

Because $x=\log M/\log2$, lower bounds for nonzero linear forms in two logarithms give an
effective irrationality measure: there exist computable $c>0$, $\mu>0$, and $Q_0$ such that
\begin{equation}
 \left|x-\frac{P}{Q}\right|>cQ^{-\mu}
 \label{eq:irrmeasure}
\end{equation}
for all $Q\ge Q_0$ unless $x=P/Q$; see, for example,
Baker--W\"ustholz and the refinements surveyed in the unified report
\cite{BakerWustholz,BakerBook}.

\begin{criterion}[Denominator-compression closure]
\label{crit:compression}
Fix $m$.  Suppose there is a sequence $r_j=e^{h_j}\in\Gamma$, $h_j\to0$, for which the reduced
denominators $Q_j$ of $R_{m,x}(r_j)$ satisfy
\begin{equation}
 C_{x,m}|h_j|^{2m}<cQ_j^{-\mu}
 \label{eq:compressionclosure}
\end{equation}
with $C_{x,m}$ an explicit remainder constant in
Theorem~\ref{thm:richardsonorder}.  Then $x\in\Z$.
\end{criterion}

\begin{proof}
If the denominators $Q_j$ are bounded along an infinite subsequence, the corresponding
approximants lie in a fixed discrete subgroup and convergence already gives $x\in\Q$.
Otherwise choose a term with $Q_j\ge Q_0$.  The upper bound
\eqref{eq:Richardsonerror}, together with \eqref{eq:compressionclosure}, contradicts
\eqref{eq:irrmeasure} unless $x=R_{m,x}(r_j)\in\Q$.  Rational descent closes the problem.
\end{proof}

A useful asymptotic version is the following.  If
\begin{equation}
 \log Q_j\le\kappa_m\log\frac1{|h_j|}+O_m(1),
 \label{eq:polyheight}
\end{equation}
then \eqref{eq:compressionclosure} follows whenever $2m>\mu\kappa_m$.  Thus arbitrary analytic
order becomes decisive if the denominator exponent $\kappa_m$ grows sublinearly in $m$.
The current unreduced bounds do not have this form: along ordinary convergents,
$\log Q_j$ is allowed to be proportional to $1/|h_j|$, not to
$\log(1/|h_j|)$.

\subsection{Why increasing \texorpdfstring{$m$}{m} does not automatically fix the scale}

One might try to choose $m=m(q)\to\infty$.  Ignoring remainder constants, the analytic gain is
roughly
\[
 |h|^{2m}\asymp q^{-2m}=\exp(-2m\log q),
\]
whereas the crude common denominator in \eqref{eq:variablemheight} has logarithm of order
$m^2q$.  Even a more economical denominator assembled from the largest term typically carries
an $mq$ scale.  Since $q\gg\log q$, the extra order does not cross the arithmetic spacing
threshold.  A proof must therefore obtain \emph{cancellation in reduced denominators}, not
merely more Taylor coefficients.

\begin{warning}
No lower bound for the reduced denominator of $R_{m,x}(r)$ is asserted here.  Large exact gcds
could radically reduce it, and finding such gcds is one of the best remaining opportunities.
The audit says only that the presently visible, unreduced denominator is too large to combine
with the analytic error.
\end{warning}

\section{New closing statements extracted from the analysis}
\label{sec:closingstatements}

The following statements are not reformulations of the conjecture by tautology; each isolates
a concrete additional arithmetic property that would be sufficient.

\begin{criterion}[Non-split extension closure]
\label{crit:nonsplit}
Let $A$ be any rational semisimple matrix with positive spectrum in $\Gamma$.  If there exists
$H\in\Mat_n(\Q)$ with nonzero class in
$\Mat_n(\Q)/[A,\Mat_n(\Q)]$ such that
$\cE_A(H)^x$ is rational, then $x\in\Z$.
\end{criterion}

This is Corollary~\ref{cor:oneselfextension}, restated as a research interface.

\begin{criterion}[Common-conductor confluence]
\label{crit:commonconductor}
Suppose there are integral $3$-smooth pairs $(a_j,b_j)$ with $b_j/a_j\to1$ such that the
reduced normalized secants
\begin{equation}
 L_j=\frac{a_j}{a_j^x}\frac{b_j^x-a_j^x}{b_j-a_j}
 \in\Q
 \label{eq:normalizedsecant}
\end{equation}
have bounded reduced denominators.  Then $x\in\Z$.
\end{criterion}

\begin{proof}
The normalized secants converge to $x$.  Bounded denominators put them in a fixed discrete
subgroup of $\Q$, so they are eventually constant and their limit is rational.
\end{proof}

\begin{criterion}[Affine order preservation]
\label{crit:affineorder}
Suppose there is a sequence $q_j\in\Gamma$, $q_j\to3$, and a fixed integer $D\ge1$ such that
\begin{equation}
 D\frac{q_j^x-N}{q_j-3}\in\Z
 \label{eq:affineD}
\end{equation}
for every $j$.  Then $x\in\Z$.
\end{criterion}

\begin{proof}
This is Proposition~\ref{prop:affinecriterion}(3).
\end{proof}

\begin{criterion}[Sublinear Richardson denominator exponent]
\label{crit:sublinear}
Assume the effective irrationality measure \eqref{eq:irrmeasure}.  Suppose that for infinitely
many orders $m$ one can find sequences $h\to0$ in $\log\Gamma$ such that
\[
 \log\den(R_{m,x}(e^h))
 \le(\kappa_m+o(1))\log\frac1{|h|},
 \qquad
 \kappa_m=o(m).
\]
Then $x\in\Z$.
\end{criterion}

\begin{proof}
Choose $m$ with $2m>\mu\kappa_m$, then apply
Criterion~\ref{crit:compression}.
\end{proof}

\begin{criterion}[Rational derivative of the $3$-smooth power homomorphism]
\label{crit:derivativehomomorphism}
Regard $\rho_x$ from \eqref{eq:rho} as a map on the dense subgroup $\Gamma\subset\R_{>0}$.
If its logarithmic derivative at one rational point exists in an arithmetic sense that forces
\begin{equation}
 q\frac{\rho_x'(q)}{\rho_x(q)}\in\Q
 \label{eq:logderivative}
\end{equation}
for one $q\in\Gamma$, then $x\in\Z$.
\end{criterion}

\begin{proof}
Analytically, $q\rho_x'(q)/\rho_x(q)=x$.  Rational descent finishes.
\end{proof}

The phrase ``in an arithmetic sense'' is deliberately not left vague in the concrete targets
above: bounded denominators, a preserved order, or a rational non-split block extension are
three exact realizations.

\section{Prioritized research program}
\label{sec:program}

\subsection{Priority A: manufacture a nonzero extension class}

The quotient theorem says exactly what must be built.  Search for a canonical module or lattice
on which the two known value maps act and whose self-extension class survives after applying
$x$.  Promising sources are:

\begin{enumerate}
\item affine semigroup actions combining dilation and translation;
\item modules of K\"ahler differentials or derivations of a suitable arithmetic semigroup
algebra;
\item non-semisimple integral reductions followed by a characteristic-zero lifting theorem;
\item exact deformation rings in which the tangent direction is canonical rather than chosen
externally.
\end{enumerate}

The failure test is immediate: compute the class in
$\Mat_n(\Q)/[A,\Mat_n(\Q)]$.  If it is zero, the construction is only a disguised similarity.

\subsection{Priority B: compute conductor defects on near-collisions}

For convergents $2^p\approx3^q$, compute
\begin{equation}
 \gcd(3^q-2^p,N^q-M^p)
 \label{eq:crossgcd}
\end{equation}
and the symmetric analogues appearing in \eqref{eq:Jintegerformula}.  The raw gap alone is not
the statistic of interest; the ratio of the gap to this gcd is.  Exact modular algorithms can
handle much larger exponents than explicit integer expansion by using modular exponentiation,
subresultants, and product trees.

A genuinely anomalous family with conductor defect polynomial in $q$ would materially change
the denominator audit.  Conversely, a theorem proving exponential conductor defect would
close off the entire confluence route and redirect effort toward extension manufacture.

\subsection{Priority C: reduced-denominator Richardson tables}

For each convergent $(p,q)$ and modest order $m$, compute the rational number
$R_{m,x}(3^q/2^p)$ symbolically in terms of $M,N$ and factor its denominator as far as possible
using only relations forced by $M=2^x$, $N=3^x$.  The aim is not numerical evaluation of a
nonexistent counterexample.  It is to derive universal gcd identities or divisibility
conditions in the polynomial ring $\Z[M,N]$ specialized to the power relation.

The coefficients \eqref{eq:cmk} are explicit and the cancellation moments are exact, making
this a good target for automated conjecture generation and formal checking.

\subsection{Priority D: adelic confluence}

A rational sequence can converge to an irrational real number because its denominators escape
through finite places.  The natural repair is adelic: seek $r_j\to1$ in the real topology while
controlling the secants in enough $p$-adic topologies that the product formula prevents escape.
The challenge is simultaneous approximation in the exponent lattice:
\begin{equation}
 a_j\log2+b_j\log3\to0
 \quad\text{and}\quad
 2^{a_j}3^{b_j}\to1\text{ in selected }\Q_p.
 \label{eq:adelicapprox}
\end{equation}
Congruence conditions on $(a_j,b_j)$ can enforce the $p$-adic part; geometry of numbers can
enforce the real part.  What is missing is a transfer theorem for the powered values
$M^{a_j}N^{b_j}$ strong enough to bound the secant denominators.

\subsection{Priority E: affine \texorpdfstring{$S$}{S}-unit rigidity}

The family $C(q)$ in Section~\ref{sec:affinecollision} asks for arithmetic regularity of the
map $q\mapsto q^x$ at the rational point $3$ when sampled on the $S$-unit group
$\Gamma$.  Classical $S$-unit equations show that exact additive patterns inside $\Gamma$ are
rare, but rarity alone is not enough.  A useful theorem would control the heights of divided
differences at an accumulation point, not merely count solutions to $u+v=1$.

\section{Lean formalization blueprint}
\label{sec:lean}

The new algebraic core is unusually suitable for formalization.  Most of it reduces to block
matrix multiplication, finite spectral decompositions, and rational linear algebra.  The
analytic interface can be isolated.

\subsection{Module 1: explicit \texorpdfstring{$2$--$3$}{2--3} blocks}

A first module can avoid general Fr\'echet derivatives entirely and formalize the concrete
formula \eqref{eq:explicit4by4}.  Suggested declarations are:

\begin{lstlisting}[style=leanstyle,caption={Schematic Lean targets for the explicit block model.}]
-- Names and exact types are schematic, not compiled source.

theorem blockRpow_diag_two_three
    (x : Real) (M N a b c d : Rat)
    (hM : Real.rpow 2 x = M) (hN : Real.rpow 3 x = N) :
    blockMatrix (diag ![2, 3]) ![![a,b],![c,d]] ^ x =
      explicitPoweredBlock x M N a b c d := by
  ...

theorem rational_blockRpow_iff_offDiag_of_irrat
    (hx : Irrational x) (hM : (2 : Real) ^ x = M)
    (hN : (3 : Real) ^ x = N) :
    EntriesRational (block23 H ^ x) <-> H 0 0 = 0 /\ H 1 1 = 0 := by
  ...

theorem integer_of_two_rpow_integer_of_nonsplitBlock
    (h2 : (2 : Real) ^ x \in Set.range Int.cast)
    (h3 : (3 : Real) ^ x \in Set.range Int.cast)
    (hblock : EntriesRational (block23 diagDirection ^ x)) :
    x \in Set.range Int.cast := by
  ...
\end{lstlisting}

The only analytic lemma needed is the power of a $2\times2$ Jordan block or, equivalently, the
finite Taylor formula for a square-zero perturbation.  The unified repository already contains
a related 
\lean{JordanFirstJet} endpoint, so this module should reuse rather than duplicate it.

\subsection{Module 2: commutator splitting}

The identities
\begin{align*}
 \Mat_n(\Q)&=\cZ(A)\oplus[A,\Mat_n(\Q)],\\
 \cE_A([A,Y])&=
 \begin{pmatrix}I&-Y\\0&I\end{pmatrix}
 (A\oplus A)
 \begin{pmatrix}I&Y\\0&I\end{pmatrix}
\end{align*}
are finite algebra.  A diagonal version should be formalized first; the spectral-projector
version can follow.

\begin{lstlisting}[style=leanstyle,caption={Schematic Lean targets for splitting and the quotient.}]
theorem offDiag_eq_commutator_of_diag_distinct
    (lambda : Fin n -> Rat)
    (h : Pairwise fun i j => lambda i != lambda j)
    (H : Matrix (Fin n) (Fin n) Rat)
    (hdiag : forall i, H i i = 0) :
    exists Y, H = diagonal lambda * Y - Y * diagonal lambda := by
  ...

theorem commutatorExtension_isSimilar
    (A Y : Matrix (Fin n) (Fin n) Rat) :
    blockExtension A (A * Y - Y * A) =
      U (-Y) * blockDiag A A * U Y := by
  ...

theorem quotient_powerDerivative_eq_smul
    ... :
    quotientMap (powerFrechet x A H) =
      x . quotientMap (AinvTimesApow * H) := by
  ...
\end{lstlisting}

\subsection{Module 3: the \texorpdfstring{$2,3,4$}{2, 3, 4} Newton chain}

The explicit matrix identity \eqref{eq:T234power} can be proved by polynomial interpolation on
three points, avoiding a general matrix-power library.  The parity endpoint is elementary.

\begin{lstlisting}[style=leanstyle,caption={Schematic Lean targets for the integral Newton chain.}]
theorem newtonChain234_rpow
    (hM : (2 : Real) ^ x = M) (hN : (3 : Real) ^ x = N) :
    rpowMatrix T234 x =
      !![M, N-M, (M^2-2*N+M)/2;
         0, N,   M^2-N;
         0, 0,   M^2] := by
  ...

theorem even_Msq_sub_twoN_add_M (M N : Int) :
    Even (M^2 - 2*N + M) := by
  omega_or_ring

theorem eigenlattice_index_T234 :
    abs (Matrix.det primitiveEigenvectors234) = 2 := by
  norm_num
\end{lstlisting}

\subsection{Module 4: Richardson coefficients}

This is pure finite combinatorics over $\Q$.

\begin{lstlisting}[style=leanstyle,caption={Schematic Lean targets for Richardson cancellation.}]
def richardsonCoeff (m k : Nat) : Rat :=
  (-1)^(k-1) * (2 * (m.factorial)^2) /
    ((m-k).factorial * (m+k).factorial)

theorem richardsonCoeffs_sum (hm : 1 <= m) :
    sum (fun k => richardsonCoeff m k) (Icc 1 m) = 1 := by
  ...

theorem richardsonCoeffs_moment
    (hj : 1 <= j) (hjm : j < m) :
    sum (fun k => richardsonCoeff m k * k^(2*j)) (Icc 1 m) = 0 := by
  ...

theorem richardsonCoeffs_firstSurviving :
    sum (fun k => richardsonCoeff m k * k^(2*m)) (Icc 1 m) =
      (-1)^(m-1) * (m.factorial)^2 := by
  ...
\end{lstlisting}

The cleanest proof uses Lagrange interpolation at the nodes $k^2$.  If existing Mathlib
interpolation infrastructure is awkward, a direct rational-function identity after clearing
factorials is likely shorter.

\subsection{Module 5: analytic remainder interface}

The analytic theorem can be stated independently of the arithmetic application:

\begin{lstlisting}[style=leanstyle,caption={Schematic analytic interface.}]
def symmetricSecant (x h : Real) : Real := sinh (x*h) / sinh h

theorem symmetricSecant_evenAnalytic :
    HasFPowerSeriesAt (fun h => symmetricSecant x h) evenSeries 0 := by
  ...

theorem richardsonSecant_isLittleO
    (m : Nat) (hm : 1 <= m) :
    (fun h => richardsonCombination m x h - x) =o[nhds 0]
      (fun h => h^(2*m-1)) := by
  ...
\end{lstlisting}

For the arithmetic report, a weaker explicit bound on a compact interval is sufficient and may
be easier to formalize than the exact power series.

\subsection{Dependency order and trust boundary}

A practical build order is:
\begin{enumerate}
\item Richardson coefficient identities;
\item explicit matrix and similarity identities;
\item the $2$--$3$ rationality test;
\item the $2,3,4$ chain and parity;
\item diagonal finite-dimensional derivative formula;
\item general spectral-projector formulation;
\item analytic Richardson remainder;
\item final conditional closing criteria.
\end{enumerate}

The first four stages require no sophisticated real matrix power API.  They can therefore be
checked quickly and provide a reliable algebraic core even if the general holomorphic
functional-calculus layer takes longer.

\section{Conclusions}
\label{sec:conclusion}

The strongest new result is not another sufficient condition of the form ``if one somehow
knows a derivative, then $x$ is rational.''  It is a complete theorem about \emph{which}
first-order extensions can remain rational under a hypothetical irrational exponent:
precisely those that were already split.  This gives an exact quotient on which power acts as
$x$ times a rational automorphism.  Any successful finite-dimensional route must put a nonzero
rational vector into that quotient.

The toric constructions exhaust a large and tempting class of candidates and provably miss the
quotient.  Integral lattice gluing and repeated tensor eigenvalues do not help by themselves;
the explicit matrix $T$ and its tensor square demonstrate the failure concretely.  Singular
conjugation does recover the Jordan jet in the limit, which explains why secant methods feel so
close, but the conjugators and normalizing denominators become unbounded.

The symmetric/Richardson hierarchy is a real quantitative improvement: it converts the prior
first-order rational secants into rational approximations with arbitrary fixed even contact
order.  It also makes the remaining obstruction more unforgivingly precise.  The missing lemma
must compress reduced denominators, preserve a nontrivial conductor, or supply an adelic
compactness mechanism.  More Taylor expansion alone will not suffice.

The most promising next move is therefore two-pronged:
\begin{enumerate}
\item algebraically, search for a canonical non-split extension arising from the additive
relation $3=2+1$ or from a derivation of an arithmetic semigroup algebra;
\item computationally, measure the common conductor and reduced Richardson denominators on very
large $2^p$--$3^q$ near-collisions, looking for exact gcd structure rather than merely small
real gaps.
\end{enumerate}

These targets are narrow enough for formalization and exact computation, yet genuinely outside
the semisimple barrier already mapped by the unified report.  They do not prove the conjecture,
but they materially reduce the space in which a proof can hide.

\appendix

\section{Quantitative Richardson bounds}
\label{app:remainder}

The main text uses $O_{x,m}(h^{2m})$, which is sufficient for the structural conclusions.  For
unit-locking experiments one wants an explicit error bound.

\begin{proposition}[A Cauchy remainder bound]
\label{prop:cauchyremainder}
Fix $x\in\R$, $m\ge1$, and $0<\rho<\pi$.  Let
\begin{equation}
 F_x(z)=\begin{cases}
 \sinh(xz)/\sinh z,&z\ne0,\\
 x,&z=0,
 \end{cases}
 \label{eq:Fx}
\end{equation}
and put
\begin{equation}
 M_{x,\rho}=\max_{|z|=\rho}|F_x(z)|,
 \qquad
 C_m=\sum_{k=1}^m|c_{m,k}|.
 \label{eq:MxrCm}
\end{equation}
If $m|h|<\rho$, then
\begin{equation}
 |R_{m,x}(e^h)-x|
 \le
 C_mM_{x,\rho}
 \frac{(m|h|/\rho)^{2m}}{1-(m|h|/\rho)^2}.
 \label{eq:explicitremainder}
\end{equation}
\end{proposition}

\begin{proof}
The function $F_x$ is even and holomorphic on $|z|<\pi$ because the zero of $\sinh z$ at zero
is removable and the next zeros are at $\pm\pi i$.  Write
$F_x(z)=x+\sum_{j\ge1}a_jz^{2j}$.  Cauchy's estimate gives
$|a_j|\le M_{x,\rho}\rho^{-2j}$.  Moment cancellation removes all $j<m$, so
\begin{align*}
 |R_{m,x}(e^h)-x|
 &\le\sum_{k=1}^m|c_{m,k}|\sum_{j\ge m}|a_j||kh|^{2j}\\
 &\le C_mM_{x,\rho}\sum_{j\ge m}(m|h|/\rho)^{2j},
\end{align*}
which is the displayed geometric sum.
\end{proof}

For reference, the first coefficient vectors and their $\ell^1$ norms are:
\begin{center}
\begin{tabular}{c>{\raggedright\arraybackslash}p{0.61\textwidth}c}
\toprule
$m$ & $(c_{m,1},\dots,c_{m,m})$ & $C_m$\\
\midrule
1 & $1$ & $1$\\
2 & $4/3,-1/3$ & $5/3$\\
3 & $3/2,-3/5,1/10$ & $11/5$\\
4 & $8/5,-4/5,8/35,-1/35$ & $93/35$\\
5 & $5/3,-20/21,5/14,-5/63,1/126$ & $193/63$\\
6 & $12/7,-15/14,10/21,-1/7,2/77,-1/462$ & $793/231$\\
\bottomrule
\end{tabular}
\end{center}

The growth of $C_m$ is modest compared with the factorial in the first surviving moment, but
this does not change the denominator-scale diagnosis: the dominant arithmetic cost comes from
the sampled rational values, not from the Richardson coefficients.

\section{Further structural consequences of the extension theorem}
\label{app:extensions}

\subsection{Multiple eigenvalues}

Theorem~\ref{thm:rationalitylocus} allows each spectral projector $P_i$ to have arbitrary rank.
On the block $P_i\Mat_n(\Q)P_i$, the derivative is multiplication by the scalar
$x\lambda_i^{x-1}$.  Thus every nonzero rational matrix coefficient within that block detects
$x$.  The quotient action is not one-dimensional; it is a direct sum of full matrix spaces,
but the same scalar $x$ appears on every summand after rational rescaling.

In particular, extra multiplicity does not weaken the detector.  What weakens it is
semisimplicity: a torus construction supplies no nonzero class in any of those matrix spaces.

\subsection{Higher triangular blocks}

For an invertible $A$ and a nilpotent $N$ commuting with $A$ and satisfying $N^{r+1}=0$, one has the terminating
formula
\begin{equation}
 (A+N)^x
 =\sum_{j=0}^r\binom{x}{j}A^{x-j}N^j.
 \label{eq:higherjet}
\end{equation}
If $A^x$ is rational and the coefficient of any nonzero $N^j$ can be shown rational, then one
obtains arithmetic information about the falling factorial
$x(x-1)\cdots(x-j+1)$.  The first coefficient already suffices, so higher jets are not needed
for logical closure.  They may nevertheless be easier to manufacture as limits of symmetric
powers, as Corollary~\ref{cor:sympowerjet} shows.

\subsection{Derivations of the value algebra}

Let $R$ be the rational algebra generated by a semisimple $A$.  A self-extension direction
commuting with $A$ can be interpreted as adjoining a square-zero element $\varepsilon H$ and
asking for a lift of the map
\[
 A\longmapsto A^x
\]
to $R[\varepsilon]/(\varepsilon^2)$.  The lift necessarily sends
\[
 A+\varepsilon H
 \longmapsto
 A^x+\varepsilon xA^{x-1}H.
\]
Thus the missing datum is a derivation-level lift of the value homomorphism.  Semisimple torus
representations determine the map on reduced algebras but provide no canonical square-zero
thickening.  This algebraic language may be a more natural starting point than matrices for
constructing the required non-split class.

\section{Symbolic verification script}
\label{app:symbolic}

The following short SymPy script independently checks the identities that are most vulnerable
to transcription errors: the centralizer/commutator block formula in the $2$--$3$ model, the
$2,3,4$ Newton-chain matrix, the first terms of the symmetric-secant expansion, the
Richardson moment identities through order $10$, and the sign in the split-to-Jordan
degeneration.

\begin{lstlisting}[style=pythonstyle,caption={Exact symbolic audit used for this report.}]
from __future__ import annotations

from math import factorial
import sympy as s

# Symbols.
x, h = s.symbols("x h")
M, N = s.symbols("M N")
a, b, c, d = s.symbols("a b c d")

# 1. The 2--3 Frechet derivative formula from spectral blocks.
A = s.diag(2, 3)
H = s.Matrix([[a, b], [c, d]])
P2 = s.diag(1, 0)
P3 = s.diag(0, 1)
D = (
    x * M / 2 * (P2 * H * P2)
    + x * N / 3 * (P3 * H * P3)
    + (N - M) * (P2 * H * P3 + P3 * H * P2)
)
assert D == s.Matrix([[x*M*a/2, (N-M)*b],
                      [(N-M)*c, x*N*d/3]])

# 2. The 2,3,4 interpolation polynomial and its value at T.
z = s.symbols("z")
p = s.interpolate([(2, M), (3, N), (4, M**2)], z)
T = s.Matrix([[2, 1, 0], [0, 3, 1], [0, 0, 4]])
I = s.eye(3)
pT = s.expand(p).coeff(z, 0) * I
pT += s.expand(p).coeff(z, 1) * T
pT += s.expand(p).coeff(z, 2) * (T*T)
expected = s.Matrix([
    [M, N-M, (M**2-2*N+M)/2],
    [0, N, M**2-N],
    [0, 0, M**2],
])
assert s.simplify(pT - expected) == s.zeros(3)

# 3. Symmetric-secant Taylor series.
series = s.series(s.sinh(x*h) / s.sinh(h), h, 0, 8)
expected_series = (
    x + x*(x**2-1)*h**2/6
    + x*(3*x**4-10*x**2+7)*h**4/360
    + x*(3*x**6-21*x**4+49*x**2-31)*h**6/15120
)
assert s.expand(series.removeO() - expected_series) == 0

# 4. Richardson coefficients and exact moments.
def coeff(m: int, k: int) -> s.Rational:
    return s.Rational(
        (-1)**(k-1) * 2 * factorial(m)**2,
        factorial(m-k) * factorial(m+k),
    )

for m in range(1, 11):
    cs = [coeff(m, k) for k in range(1, m+1)]
    assert sum(cs) == 1
    for j in range(1, m):
        assert sum(cs[k-1] * k**(2*j)
                   for k in range(1, m+1)) == 0
    assert sum(cs[k-1] * k**(2*m)
               for k in range(1, m+1)) == (-1)**(m-1) * factorial(m)**2

# 5. Correct sign in the split-to-Jordan conjugation.
a0, r, k = s.symbols("a0 r k", nonzero=True)
b0 = a0*r
S = s.Matrix([[1, -k], [0, 1]])
B = s.simplify(S.inv() * s.diag(a0, b0) * S)
expected_B = s.Matrix([[a0, k*(b0-a0)], [0, b0]])
assert s.simplify(B - expected_B) == s.zeros(2)

print("all exact audits passed")
\end{lstlisting}

The script uses exact rational and symbolic arithmetic.  It is a transcription check, not a
substitute for the proofs in the main text and not a proof-assistant certificate.

\section{Status and dependency matrix}
\label{app:status}

\begin{longtable}{>{\raggedright\arraybackslash}p{0.34\textwidth}
                  >{\raggedright\arraybackslash}p{0.18\textwidth}
                  >{\raggedright\arraybackslash}p{0.38\textwidth}}
\caption{Status of the principal statements.}\label{tab:status}\\
\toprule
Statement & Status & Dependencies and caveats\\
\midrule
\endfirsthead
\toprule
Statement & Status & Dependencies and caveats\\
\midrule
\endhead
Block derivative identity, Theorem~\ref{thm:blockderivative} & \paperproof{} & Standard
holomorphic functional calculus; contour proof included\\
\addlinespace
Semisimple divided-difference formula, Theorem~\ref{thm:semisimplederivative} & \paperproof{} &
Finite spectral decomposition\\
\addlinespace
Exact rationality locus, Theorem~\ref{thm:rationalitylocus} & \paperproof{} & Requires rational
spectral projectors and rational powered eigenvalues; no transcendence conjecture\\
\addlinespace
Extension quotient action, Theorem~\ref{thm:quotientaction} & \paperproof{} & Basis-free
consequence of the rationality-locus decomposition\\
\addlinespace
Toric-functor no-go, Theorem~\ref{thm:toricnogo} & \paperproof{} & Standard semisimplicity of
split-torus representations in characteristic zero\\
\addlinespace
$2,3,4$ Newton chain, Theorem~\ref{thm:T234power} & \paperproof{} & Exact interpolation and a
parity check; symbolically audited\\
\addlinespace
Split-to-Jordan degeneration, Theorem~\ref{thm:splitdegeneration} & \paperproof{} & Uses
unbounded integral conjugators; gives a limit, not arithmetic compactness\\
\addlinespace
Symmetric secant expansion, Theorem~\ref{thm:secondorder} & \paperproof{} & Elementary power
series; symbolically audited\\
\addlinespace
Richardson moment identities and $O(h^{2m})$ theorem & \paperproof{} & Lagrange interpolation
at nodes $k^2$; symbolically audited through $m=10$\\
\addlinespace
Explicit Cauchy remainder, Proposition~\ref{prop:cauchyremainder} & \paperproof{} & Valid for
$m|h|<\rho<\pi$\\
\addlinespace
Denominator-compression and locking criteria & \paperproof{} & Conditional statements; the
required denominator estimates remain open\\
\addlinespace
Full \AEname{} conjecture & \openstatus{} & Not proved here or in the attached manuscript\\
\bottomrule
\end{longtable}

\begin{thebibliography}{99}

\bibitem{UnifiedReport2026}
The ProveIt project,
\newblock \emph{The Alaoglu--Erd\H{o}s Integer-Exponent Conjecture: Machine-Checked Partial
Results, Exact Conditional Structure, Determinant Methods, and the Remaining Barrier},
\newblock unified working manuscript with twenty-seven continuation reports, August 2026;
attached source file \texttt{alaoglu-erdos-unified-report.tex(3).gz}.

\bibitem{AE1944}
L.~Alaoglu and P.~Erd\H{o}s,
\newblock \emph{On highly composite and similar numbers},
\newblock Transactions of the American Mathematical Society \textbf{56} (1944), 448--469.
\newblock DOI:
\href{https://doi.org/10.1090/S0002-9947-1944-0011087-2}{10.1090/S0002-9947-1944-0011087-2}.

\bibitem{BakerBook}
A.~Baker,
\newblock \emph{Transcendental Number Theory},
\newblock Cambridge University Press, 1975.

\bibitem{BakerWustholz}
A.~Baker and G.~W\"ustholz,
\newblock \emph{Logarithmic forms and group varieties},
\newblock Journal f\"ur die reine und angewandte Mathematik \textbf{442} (1993), 19--62.
\newblock DOI:
\href{https://doi.org/10.1515/crll.1993.442.19}{10.1515/crll.1993.442.19}.

\bibitem{DasguptaKakde}
S.~Dasgupta and M.~Kakde,
\newblock \emph{Ranks of matrices of logarithms of algebraic numbers II: The Matrix
Coefficient Conjecture},
\newblock Advances in Mathematics \textbf{487} (2026), Article 110753;
arXiv:2408.08178.
\newblock DOI:
\href{https://doi.org/10.1016/j.aim.2025.110753}{10.1016/j.aim.2025.110753}.

\bibitem{Higham}
N.~J.~Higham,
\newblock \emph{Functions of Matrices: Theory and Computation},
\newblock SIAM, 2008.
\newblock DOI: \href{https://doi.org/10.1137/1.9780898717778}{10.1137/1.9780898717778}.

\bibitem{Kato}
T.~Kato,
\newblock \emph{Perturbation Theory for Linear Operators},
\newblock second edition, Springer, 1976; corrected printing, 1995.

\bibitem{Richardson1911}
L.~F.~Richardson,
\newblock \emph{The approximate arithmetical solution by finite differences of physical
problems involving differential equations, with an application to the stresses in a masonry
dam},
\newblock Philosophical Transactions of the Royal Society A \textbf{210} (1911), 307--357.

\bibitem{AnthropicFable}
Anthropic,
\newblock \emph{Claude Fable 5}, official model page and announcement, June--July 2026.
\newblock \url{https://www.anthropic.com/claude/fable}.

\bibitem{AnthropicRedeploy}
Anthropic,
\newblock \emph{Redeploying Fable 5}, June 30, 2026; access restored July 1, 2026.
\newblock \url{https://www.anthropic.com/news/redeploying-fable-5}.

\bibitem{GaoJacobian}
S.~Gao,
\newblock \emph{Counterexamples to the Jacobian conjecture in dimensions greater than two},
\newblock arXiv:2608.00222, July 31, 2026.
\newblock \url{https://arxiv.org/abs/2608.00222}.

\bibitem{SilvermanMO}
J.~H.~Silverman,
\newblock update to \emph{What heuristic evidence is there concerning the unboundedness or
boundedness of Mordell--Weil ranks of elliptic curves over $\Q$?},
\newblock MathOverflow, August 2026.
\newblock \url{https://mathoverflow.net/questions/38356/}.

\bibitem{ICARMRank}
Institute for Computer-Aided Reasoning in Mathematics,
\newblock \emph{Elliptic Curve Rank Leaderboard}, accessed August 22, 2026.
\newblock \url{https://elliptic-rank.icarm.cloud/}.

\bibitem{VibeMathedHadamard}
VibeMathed contributors,
\newblock \emph{Hadamard Matrix of Order 668}, exact AI-assisted construction record and
verification note, August 2026.
\newblock \url{https://vibemathed.com/}.

\bibitem{OEISHadamard}
OEIS Foundation,
\newblock \emph{A007299: orders for which a Hadamard matrix is not yet known},
\newblock update recording the twelve 2026 constructions below $2000$, accessed August 22,
2026.
\newblock \url{https://oeis.org/A007299}.

\end{thebibliography}

\end{document}
